% Problems, Chapter 2: Determinants and Matrices
% Arfken, Weber & Harris, Mathematical Methods for Physicists, 7th ed.
% Requires: amsmath, amssymb, tikz

\section*{Problems}

\Needspace{9\baselineskip}
\subsection*{2.1 Determinants}

\Needspace{5\baselineskip}
\par\bigskip
\noindent\textbf{2.1.1}\quad Evaluate the following determinants:
\begin{enumerate}
\item[(a)] $\begin{vmatrix}1&0&1\\0&1&0\\1&0&0\end{vmatrix}$,\quad
\item[(b)] $\begin{vmatrix}1&2&0\\3&1&2\\0&3&1\end{vmatrix}$,\quad
\item[(c)] $\dfrac{1}{\sqrt{2}}\begin{vmatrix}
  0&\sqrt{3}&0&0\\ \sqrt{3}&0&2&0\\ 0&2&0&\sqrt{3}\\
  0&0&\sqrt{3}&0\end{vmatrix}$.
\end{enumerate}
\solution
Expand along a row or column containing as many zeros as possible.
\begin{enumerate}
\item[(a)] Only the middle entry of row 2 contributes:
\[
D_a=1(-1)^{2+2}
\begin{vmatrix}1&1\\1&0\end{vmatrix}
=1(1\cdot0-1\cdot1)=\boxed{-1}.
\]
\item[(b)] Expansion along row 1 gives
\[
D_b=1\begin{vmatrix}1&2\\3&1\end{vmatrix}
-2\begin{vmatrix}3&2\\0&1\end{vmatrix}
 =1(1-6)-2(3-0)=\boxed{-11}.
\]
\item[(c)] Write $s=\sqrt3$. The factor $1/\sqrt2$ in the question
multiplies the determinant once; it is not a factor on each row.
Expand along row 1, whose only nonzero entry is $a_{12}=s$ with sign
$(-1)^{1+2}=-1$, and then along column 1 of the resulting minor, whose
only nonzero entry is $s$ in position $(1,1)$:
\[
\begin{aligned}
D_c
&=-\frac{s}{\sqrt2}
\begin{vmatrix}s&2&0\\0&0&s\\0&s&0\end{vmatrix}
 =-\frac{s^2}{\sqrt2}
\begin{vmatrix}0&s\\s&0\end{vmatrix}\\
&=-\frac{3}{\sqrt2}(0-3)=\boxed{\frac9{\sqrt2}}.
\end{aligned}
\]
\end{enumerate}

\Needspace{5\baselineskip}
\par\bigskip
\noindent\textbf{2.1.2}\quad Test the set of linear homogeneous equations
\[
x+3y+3z=0,\quad x-y+z=0,\quad 2x+y+3z=0
\]
to see if it possesses a nontrivial solution. In any case, find a solution to
this equation set.
\solution
The coefficient determinant, expanded along row 1, is
\[
\begin{aligned}
D&=\begin{vmatrix}1&3&3\\1&-1&1\\2&1&3\end{vmatrix}\\
 &=1(-3-1)-3(3-2)+3(1+2)=-4-3+9=2.
\end{aligned}
\]
Since $D\ne0$, only the trivial solution is possible. Directly,
the second equation gives $x=y-z$. The first and third then become
\[
(y-z)+3y+3z=4y+2z=0,\qquad
2(y-z)+y+3z=3y+z=0.
\]
Thus $z=-2y$ and $z=-3y$, so $y=0$, $z=0$, and $x=0$:
$\boxed{(x,y,z)=(0,0,0)}$.
Substitution into each original equation gives zero.

\Needspace{5\baselineskip}
\par\bigskip
\noindent\textbf{2.1.3}\quad Given the pair of equations
\[
x+2y=3,\quad 2x+4y=6,
\]
\begin{enumerate}
\item[(a)] Show that the determinant of the coefficients vanishes.
\item[(b)] Show that the numerator determinants, see Eq.~(2.18), also vanish.
\item[(c)] Find at least two solutions.
\end{enumerate}
\solution
\textbf{(a), (b)} The coefficient determinant $D$ and the numerator
determinants $D_x$, $D_y$ of Cramer's rule are
\[
D=\begin{vmatrix}1&2\\2&4\end{vmatrix}=4-4=0,\quad
D_x=\begin{vmatrix}3&2\\6&4\end{vmatrix}=12-12=0,\quad
D_y=\begin{vmatrix}1&3\\2&6\end{vmatrix}=6-6=0.
\]
\textbf{(c)} The quotients $D_x/D,D_y/D$ are undefined and cannot
determine a unique solution. The second equation is twice the first,
so put $y=t$. Then $x=3-2t$, and every real solution is
\[
\boxed{(x,y)=(3-2t,t),\quad t\in\mathbb R}.
\]
(Complex $t$ is allowed if complex solutions are sought.)
For example $t=0$ gives $(3,0)$, and $t=1$ gives $(1,1)$.
Both satisfy $x+2y=3$ and hence $2x+4y=6$.

\Needspace{5\baselineskip}
\par\bigskip
\noindent\textbf{2.1.4}\quad If $C_{ij}$ is the cofactor of element $a_{ij}$,
formed by striking out the $i$th row and $j$th column and including a sign
$(-1)^{i+j}$, show that
\begin{enumerate}
\item[(a)] $\displaystyle\sum_{i}a_{ij}C_{ij}=\sum_{i}a_{ji}C_{ji}=|\mathsf{A}|$,
  where $|\mathsf{A}|$ is the determinant with the elements $a_{ij}$,
\item[(b)] $\displaystyle\sum_{i}a_{ij}C_{ik}=\sum_{i}a_{ji}C_{ki}=0$,
  $j\neq k$.
\end{enumerate}
\solution
Let $A_{\widehat i,\widehat j}$ denote $A$ with row $i$ and column $j$
struck out, so that $C_{ij}=(-1)^{i+j}\det A_{\widehat i,\widehat j}$.
For fixed $j$, Laplace expansion down column $j$ is
$\det A=\sum_i a_{ij}(-1)^{i+j}\det A_{\widehat i,\widehat j}
=\sum_i a_{ij}C_{ij}$.
Expansion across row $j$ similarly gives
$\det A=\sum_i a_{ji}C_{ji}$. This proves (a), without assuming
invertibility.

For (b), replace column $k$ of $A$ by its column $j$, with $j\ne k$,
and call the result $B$. Columns $j$ and $k$ of $B$ are identical, so
$\det B=0$. Expanding down column $k$, each minor is unchanged
from $A$ because that column is deleted; its cofactor is $C_{ik}$.
The entry being multiplied by that cofactor is now $a_{ij}$. Therefore
\[
0=\det B=\sum_i a_{ij}C_{ik}.
\]
Replacing row $k$ by row $j$ and expanding across row $k$ gives,
by the same argument, $0=\sum_i a_{ji}C_{ki}$.

\Needspace{5\baselineskip}
\par\bigskip
\noindent\textbf{2.1.5}\quad A determinant with all elements of order unity
may be surprisingly small. The Hilbert determinant
$H_{ij}=(i+j-1)^{-1}$, $i,j=1,2,\ldots,n$ is notorious for its small values.
\begin{enumerate}
\item[(a)] Calculate the value of the Hilbert determinants of order $n$ for
  $n=1$, $2$, and $3$.
\item[(b)] If an appropriate subroutine is available, find the Hilbert
  determinants of order $n$ for $n=4$, $5$, and $6$.
\end{enumerate}
\par\nobreak\hfill\emph{ANS.}\enspace
$\begin{array}{cr}
n&\mathrm{Det}(H_{n})\\
1&1.\\
2&8.33333\times10^{-2}\\
3&4.62963\times10^{-4}\\
4&1.65344\times10^{-7}\\
5&3.74930\times10^{-12}\\
6&5.36730\times10^{-18}.
\end{array}$
\solution
\textbf{(a)} Use exact fractions so that small determinants are not lost
by subtracting rounded numbers. The first three determinants ($D_3$
expanded along row 1) are
\[
D_1=1,\qquad D_2=1\cdot\frac13-\frac12\cdot\frac12=\frac1{12},
\]
\[
\begin{aligned}
D_3
&=\left(\frac1{15}-\frac1{16}\right)
-\frac12\left(\frac1{10}-\frac1{12}\right)
+\frac13\left(\frac18-\frac19\right)\\
&=\frac1{240}-\frac1{120}+\frac1{216}
=\frac{9-18+10}{2160}=\frac1{2160}.
\end{aligned}
\]

\textbf{(b)} For the larger matrices, Gaussian elimination uses
\[
R_i\longleftarrow R_i-\frac{u_{ik}}{u_{kk}}R_k
\quad(i=k+1,\ldots,n;\ k=1,\ldots,n-1),
\]
where $u_{ik}$ is the current entry in row $i$, column $k$; row $k$ no
longer changes after step $k-1$, so $u_{kk}$ is the $k$th pivot.
Each operation preserves the determinant. For example, the first
step on $H_6$ changes an entry with $i,j\ge2$ into
\[
\frac1{i+j-1}-\frac1i\frac1j
=\frac{(i-1)(j-1)}{ij(i+j-1)}.
\]
Every later step has the same structure. Suppose that before step $k$
the entries with $i,j\ge k$ have the form $\alpha_i\beta_j/(i+j-1)$;
for $H_6$ itself $\alpha_i=\beta_j=1$. The multiplier for row $i$ is
then
\[
\frac{u_{ik}}{u_{kk}}
=\frac{\alpha_i\beta_k/(i+k-1)}{\alpha_k\beta_k/(2k-1)}
=\frac{\alpha_i}{\alpha_k}\,\frac{2k-1}{i+k-1}.
\]
Subtracting it times the row-$k$ entry $\alpha_k\beta_j/(k+j-1)$ and
using the identity (expand both products)
\[
(i+k-1)(j+k-1)-(2k-1)(i+j-1)=ij-k(i+j)+k^2=(i-k)(j-k)
\]
turns the entry in row $i$, column $j$ into
\[
\alpha_i\beta_j\left[\frac1{i+j-1}-\frac{2k-1}{(i+k-1)(j+k-1)}\right]
=\frac{\alpha_i\beta_j\,(i-k)(j-k)}{(i+k-1)(j+k-1)(i+j-1)}.
\]
The form is preserved, with $\alpha_i$ multiplied by $(i-k)/(i+k-1)$
and $\beta_j$ by $(j-k)/(j+k-1)$; for $k=1$ this is the formula above.
Hence at step $k$
\[
\alpha_k=\beta_k=\prod_{m=1}^{k-1}\frac{k-m}{k+m-1}
=\frac{[(k-1)!]^2}{(2k-2)!},\qquad
u_{kk}=\frac{\alpha_k\beta_k}{2k-1}
=\frac{[(k-1)!]^4}{(2k-2)!\,(2k-1)!},
\]
because $\prod_{m=1}^{k-1}(k-m)=(k-1)!$,
$\prod_{m=1}^{k-1}(k+m-1)=k(k+1)\cdots(2k-2)=(2k-2)!/(k-1)!$, and
$(2k-2)!\,(2k-1)=(2k-1)!$. Successive pivots therefore have the ratio
\[
\frac{u_{k+1,k+1}}{u_{kk}}
=\frac{k^4}{(2k-1)(2k)\cdot(2k)(2k+1)}
=\frac{k^2}{4(4k^2-1)}
=\frac1{12},\ \frac1{15},\ \frac9{140},\ \frac4{63},\ \frac{25}{396}
\quad(k=1,\ldots,5),
\]
which, starting from $u_{11}=1$, gives the pivots $1$, $1/12$, $1/180$,
$1/2800$, $1/44100$, $1/698544$ for $k=1,\ldots,6$. The other entries
of row $k$ of the final triangular matrix are
$u_{kj}=\alpha_k\beta_j/(k+j-1)$ for $j>k$,
with $\beta_j=\prod_{m=1}^{k-1}(j-m)/(j+m-1)$.
For completeness, the multipliers $u_{ik}/u_{kk}$ in the successive
steps are listed below. At step $k$, subtract the indicated multiple
of the current row $k$ from each listed row $i$:
\[
\begin{array}{c|c|c}
k&i&u_{ik}/u_{kk}\ \text{in the same order}\\\hline
1&2,3,4,5,6&1/2,\ 1/3,\ 1/4,\ 1/5,\ 1/6\\
2&3,4,5,6&1,\ 9/10,\ 4/5,\ 5/7\\
3&4,5,6&3/2,\ 12/7,\ 25/14\\
4&5,6&2,\ 25/9\\
5&6&5/2
\end{array}
\]
For example, the second pivot is $1/12$ and the entry below it
in row 4 is $3/40$, giving multiplier $(3/40)/(1/12)=9/10$.
These operations produce
\[
U_6=\begin{pmatrix}
1&1/2&1/3&1/4&1/5&1/6\\
0&1/12&1/12&3/40&1/15&5/84\\
0&0&1/180&1/120&1/105&5/504\\
0&0&0&1/2800&1/1400&1/1008\\
0&0&0&0&1/44100&1/17640\\
0&0&0&0&0&1/698544
\end{pmatrix}.
\]
Elimination on a leading $n\times n$ block uses the same leading
pivots, so $D_n$ is the product of the first $n$ diagonal entries.
In particular
\[
D_4=\frac{D_3}{2800},\qquad
D_5=\frac{D_4}{44100},\qquad
D_6=\frac{D_5}{698544}.
\]
The resulting exact values and six-significant-figure approximations are
\[
\boxed{\begin{array}{c|c|c}
n&D_n&\text{decimal value}\\\hline
1&1&1\\
2&1/12&8.33333\times10^{-2}\\
3&1/2160&4.62963\times10^{-4}\\
4&1/6048000&1.65344\times10^{-7}\\
5&1/266716800000&3.74930\times10^{-12}\\
6&1/186313420339200000&5.36730\times10^{-18}
\end{array}}
\]
Equivalently, $D_n=\prod_{k=1}^n[(k-1)!]^4/[(2k-2)!\,(2k-1)!]$.
As an independent check, evaluating $\det H_n$ directly in exact
rational arithmetic (for instance by cofactor expansion in a
computer-algebra system) reproduces the same six fractions.

\Needspace{5\baselineskip}
\par\bigskip
\noindent\textbf{2.1.6}\quad Prove that the determinant consisting of the
coefficients from a set of linearly dependent forms has the value zero.
\solution
Write the $n$ forms as $L_i(x)=\sum_{j=1}^n a_{ij}x_j$.
Linear dependence means that constants $c_i$, not all zero, satisfy
$\sum_i c_iL_i(x)=0$ for every $x$.
Since $\sum_i c_iL_i(x)=\sum_j\bigl(\sum_i c_ia_{ij}\bigr)x_j$, the
choice $x=e_j$ (the $j$th unit vector) gives $\sum_i c_i a_{ij}=0$ for
each $j$; that is, the rows $R_i=(a_{i1},\ldots,a_{in})$ of the
coefficient matrix $A$ satisfy $\sum_i c_iR_i=0$.
Choose $r$ with $c_r\ne0$. Then the row vectors obey
\[
R_r=-\sum_{i\ne r}\frac{c_i}{c_r}R_i.
\]
Replace $R_r$ by $R_r+\sum_{i\ne r}(c_i/c_r)R_i$.
Adding multiples of other rows leaves the determinant unchanged, but
the new row is zero. Expansion across it therefore gives
$\boxed{\det A=0}$.
The same proof applies if the forms are recorded as columns.

\Needspace{5\baselineskip}
\par\bigskip
\noindent\textbf{2.1.7}\quad Solve the following set of linear simultaneous
equations. Give the results to five decimal places.
\[
\begin{split}
1.0x_{1}+0.9x_{2}+0.8x_{3}+0.4x_{4}+0.1x_{5}&=1.0\\
0.9x_{1}+1.0x_{2}+0.8x_{3}+0.5x_{4}+0.2x_{5}+0.1x_{6}&=0.9\\
0.8x_{1}+0.8x_{2}+1.0x_{3}+0.7x_{4}+0.4x_{5}+0.2x_{6}&=0.8\\
0.4x_{1}+0.5x_{2}+0.7x_{3}+1.0x_{4}+0.6x_{5}+0.3x_{6}&=0.7\\
0.1x_{1}+0.2x_{2}+0.4x_{3}+0.6x_{4}+1.0x_{5}+0.5x_{6}&=0.6\\
0.1x_{2}+0.2x_{3}+0.3x_{4}+0.5x_{5}+1.0x_{6}&=0.5.
\end{split}
\]
\emph{Note.}~These equations may also be solved by matrix inversion, as
discussed in Section 2.2.
\solution
Multiply every equation by $10$ and put them in the order $5,6,3,4,1,2$:
\[
\left[\begin{array}{rrrrrr|r}
1&2&4&6&10&5&6\\
0&1&2&3&5&10&5\\
8&8&10&7&4&2&8\\
4&5&7&10&6&3&7\\
10&9&8&4&1&0&10\\
9&10&8&5&2&1&9
\end{array}\right].
\]
Here $R_i$ refers to this new ordering. Apply
$R_3-8R_1$, $R_4-4R_1$, $R_5-10R_1$, and $R_6-9R_1$ to
rows $3,4,5,6$. The last four rows, omitting the zero first column, become
\[
\left[\begin{array}{rrrrr|r}
-8&-22&-41&-76&-38&-40\\
-3&-9&-14&-34&-17&-17\\
-11&-32&-56&-99&-50&-50\\
-8&-28&-49&-88&-44&-45
\end{array}\right].
\]
Next add $8R_2,3R_2,11R_2,8R_2$ to rows $3,4,5,6$, respectively,
which clears the $x_2$ column.
In columns $x_3,x_4,x_5,x_6$ this gives
\[
\left[\begin{array}{rrrr|r}
-6&-17&-36&42&0\\
-3&-5&-19&13&-2\\
-10&-23&-44&60&5\\
-12&-25&-48&36&-5
\end{array}\right].
\]
To eliminate $x_3$, replace
\[
R_4\leftarrow-2R_4+R_3,\quad
R_5\leftarrow-3R_5+5R_3,\quad
R_6\leftarrow-R_6+2R_3.
\]
Rows $4,5,6$, in columns $x_4,x_5,x_6$, are now
\[
\left[\begin{array}{rrr|r}
-7&2&16&4\\
-16&-48&30&-15\\
-9&-24&48&5
\end{array}\right].
\]
Eliminate $x_4$ by
$R_5\leftarrow-7R_5+16R_4$ and
$R_6\leftarrow-7R_6+9R_4$. The final two equations are
\[
368x_5+46x_6=169,\qquad186x_5-192x_6=1.
\]
Finally $R_6\leftarrow(368R_6-186R_5)/14$ gives
$-5658x_6=-2219$, because
$368(-192)-186\cdot46=-79212=14(-5658)$ and
$368\cdot1-186\cdot169=-31066=14(-2219)$.
Hence $x_6=2219/5658$. To back-substitute without fractions, write
$x_r=y_r/5658$ and use, in turn, $368x_5+46x_6=169$,
row 4 ($-7x_4+2x_5+16x_6=4$), row 3
($-6x_3-17x_4-36x_5+42x_6=0$), and rows 2 and 1 of the reordered
matrix:
\[
\begin{aligned}
y_6&=2219,\\
y_5&=\frac{169\cdot5658-46y_6}{368}
=\frac{956202-102074}{368}=2321,\\
y_4&=\frac{2y_5+16y_6-4\cdot5658}{7}
=\frac{4642+35504-22632}{7}=2502,\\
y_3&=\frac{-17y_4-36y_5+42y_6}{6}
=\frac{-42534-83556+93198}{6}=-5482,\\
y_2&=5\cdot5658-2y_3-3y_4-5y_5-10y_6\\
&=28290+10964-7506-11605-22190=-2047,\\
y_1&=6\cdot5658-2y_2-4y_3-6y_4-10y_5-5y_6\\
&=33948+4094+21928-15012-23210-11095=10653.
\end{aligned}
\]
Since $5658=2\cdot3\cdot23\cdot41$ and the common factors of $5658$ with
$y_1,\ldots,y_6$ are $3$, $23$, $2$, $6$, $1$, $1$, the fractions
$y_r/5658$ reduce to
\[
x=\Bigl(\frac{3551}{1886},-\frac{89}{246},-\frac{2741}{2829},
\frac{417}{943},\frac{2321}{5658},\frac{2219}{5658}\Bigr)^T.
\]
Thus, to five decimal places,
\[
\boxed{x=(1.88282,-0.36179,-0.96889,0.44221,0.41022,0.39219)^T}.
\]
For an exact check, let $B$ be ten times the original coefficient
matrix and multiply it into the vector of numerators $y_r$:
\[
\begin{aligned}
B(10653,-2047,-5482,2502,2321,2219)^T
&=(56580,50922,45264,39606,33948,28290)^T\\
&=5658(10,9,8,7,6,5)^T,
\end{aligned}
\]
verifying all six equations without rounding.

\Needspace{5\baselineskip}
\par\bigskip
\noindent\textbf{2.1.8}\quad Show that (in 3-D space)
\begin{enumerate}
\item[(a)] $\displaystyle\sum_{i}\delta_{ii}=3$,
\item[(b)] $\displaystyle\sum_{ij}\delta_{ij}\varepsilon_{ijk}=0$,
\item[(c)] $\displaystyle\sum_{pq}\varepsilon_{ipq}\varepsilon_{jpq}
  =2\delta_{ij}$,
\item[(d)] $\displaystyle\sum_{ijk}\varepsilon_{ijk}\varepsilon_{ijk}=6$.
\end{enumerate}
\emph{Note.}~The symbol $\delta_{ij}$ is the Kronecker delta, defined in
Eq.~(1.164), and $\varepsilon_{ijk}$ is the Levi-Civita symbol, Eq.~(2.8).
\solution
\begin{enumerate}
\item[(a)] Since $\delta_{11}=\delta_{22}=\delta_{33}=1$,
$\sum_i\delta_{ii}=1+1+1=3$.
\item[(b)] The factor $\delta_{ij}$ eliminates every term with $i\ne j$,
so $\sum_{ij}\delta_{ij}\varepsilon_{ijk}=\sum_i\varepsilon_{iik}=0$.
Every remaining Levi-Civita symbol has a repeated index.
\item[(c)] If $i=j$, the only nonzero terms have $(p,q)$ equal to the
two orderings of the other two indices. Their squares are both $1$, so
the sum is $2$. If $i\ne j$, a nonzero $\varepsilon_{ipq}$ requires
$\{i,p,q\}=\{1,2,3\}$, so $j\ne i$ equals $p$ or $q$. For the same
$p,q$, $\varepsilon_{jpq}$
then has a repeated index and vanishes. Thus the sum is $2$ for $i=j$
and $0$ otherwise: $\sum_{pq}\varepsilon_{ipq}\varepsilon_{jpq}=2\delta_{ij}$.
\item[(d)] There are $3!=6$ permutations of $(1,2,3)$.
For each, $\varepsilon_{ijk}^2=1$, and every other term vanishes.
Thus $\sum_{ijk}\varepsilon_{ijk}\varepsilon_{ijk}=6$.
Equivalently, set $j=i$ in (c) and sum over $i$:
$\sum_i2\delta_{ii}=6$ by (a).
\end{enumerate}

\Needspace{5\baselineskip}
\par\bigskip
\noindent\textbf{2.1.9}\quad Show that (in 3-D space)
\[
\sum_{k}\varepsilon_{ijk}\varepsilon_{pqk}
=\delta_{ip}\delta_{jq}-\delta_{iq}\delta_{jp}.
\]
\emph{Note.}~See Exercise 2.1.8 for definitions of $\delta_{ij}$ and
$\varepsilon_{ijk}$.
\solution
\erratum{The textbook identity is correct. In the Instructor's Manual
the final line has
$\delta_{ip}\delta_{jp}$ in place of $\delta_{ip}\delta_{jq}$.
For $i=p=1$, $j=q=2$, the left side is $1$, whereas that misprinted
right side is $0$.}
If $i=j$ or $p=q$, a Levi-Civita factor on the left vanishes for
every $k$; the two delta products on the right are identical and cancel.
Now take $i\ne j$ and $p\ne q$. A nonzero product requires a $k$
different from both $i,j$ and from both $p,q$. In three dimensions
$k$ is then the one index missing from $\{i,j\}$ and also the one
missing from $\{p,q\}$, so the unordered pairs $\{i,j\}$ and
$\{p,q\}$ must agree. If they differ, the left side is zero, and so
is the right side, since each delta product on the right requires
$(p,q)=(i,j)$ or $(p,q)=(j,i)$. If $(p,q)=(i,j)$, the single
allowed $k$ contributes $\varepsilon_{ijk}^2=1$, and the right side is
$1-0$. If $(p,q)=(j,i)$, antisymmetry gives
$\varepsilon_{ijk}\varepsilon_{jik}=-1$, and the right side is $0-1$.
These cases exhaust all indices, proving
\[
\boxed{\sum_k\varepsilon_{ijk}\varepsilon_{pqk}
=\delta_{ip}\delta_{jq}-\delta_{iq}\delta_{jp}}.
\]

\Needspace{9\baselineskip}
\subsection*{2.2 Matrices}

\Needspace{5\baselineskip}
\par\bigskip
\noindent\textbf{2.2.1}\quad Show that matrix multiplication is associative,
$(\mathsf{AB})\mathsf{C}=\mathsf{A}(\mathsf{BC})$.
\solution
Let $A,B,C$ have sizes $m\times n$, $n\times p$, and $p\times q$.
For every allowed $i,j$,
\[
\begin{aligned}
((AB)C)_{ij}
&=\sum_{k=1}^p(AB)_{ik}c_{kj}
=\sum_{k=1}^p\sum_{\ell=1}^n a_{i\ell}b_{\ell k}c_{kj}\\
&=\sum_{\ell=1}^n a_{i\ell}
       \left(\sum_{k=1}^p b_{\ell k}c_{kj}\right)
=\sum_{\ell=1}^n a_{i\ell}(BC)_{\ell j}
=(A(BC))_{ij}.
\end{aligned}
\]
The sums are finite, so their order can be interchanged. Equality of
every entry proves $(AB)C=A(BC)$.

\Needspace{5\baselineskip}
\par\bigskip
\noindent\textbf{2.2.2}\quad Show that
\[
(\mathsf{A}+\mathsf{B})(\mathsf{A}-\mathsf{B})=\mathsf{A}^{2}-\mathsf{B}^{2}
\]
if and only if $\mathsf{A}$ and $\mathsf{B}$ commute,
\[
[\mathsf{A},\mathsf{B}]=0.
\]
\solution
Distribute while preserving factor order:
\[
(A+B)(A-B)=A^2-AB+BA-B^2=A^2-B^2-[A,B].
\]
If $[A,B]=0$, the required equality follows. Conversely, if that
equality holds, subtract $A^2-B^2$ from both sides to get
$-[A,B]=0$. This proves both directions.

\Needspace{5\baselineskip}
\par\bigskip
\noindent\textbf{2.2.3}\quad
\begin{enumerate}
\item[(a)] Complex numbers, $a+ib$, with $a$ and $b$ real, may be represented
  by (or are isomorphic with) $2\times2$ matrices:
  \[
  a+ib\longleftrightarrow\begin{pmatrix}a&b\\-b&a\end{pmatrix}.
  \]
  Show that this matrix representation is valid for (i) addition and (ii)
  multiplication.
\item[(b)] Find the matrix corresponding to $(a+ib)^{-1}$,
  \erratum{assuming $a+ib\ne0$}.
\end{enumerate}
\solution
\textbf{(a)} Define $F(a+ib)=\begin{pmatrix}a&b\\-b&a\end{pmatrix}$.
(i) For real $a,b,c,d$,
\[
F(a+ib)+F(c+id)=
\begin{pmatrix}a+c&b+d\\-(b+d)&a+c\end{pmatrix}
=F((a+c)+i(b+d)).
\]
(ii) Multiplication gives
\[
\begin{pmatrix}a&b\\-b&a\end{pmatrix}
\begin{pmatrix}c&d\\-d&c\end{pmatrix}
=\begin{pmatrix}ac-bd&ad+bc\\-(ad+bc)&ac-bd\end{pmatrix},
\]
which is $F((a+ib)(c+id))$ because
$(a+ib)(c+id)=(ac-bd)+i(ad+bc)$.
The representation is one-to-one, since its first row recovers $a,b$,
and $F(1)=I_2$.

\textbf{(b)} For $a+ib\ne0$, multiplying numerator and denominator by
$a-ib$ gives $(a+ib)^{-1}=x+iy$ with $x=a/(a^2+b^2)$,
$y=-b/(a^2+b^2)$, so
\[
\boxed{F((a+ib)^{-1})=\begin{pmatrix}x&y\\-y&x\end{pmatrix}
=\frac1{a^2+b^2}
\begin{pmatrix}a&-b\\b&a\end{pmatrix}}.
\]
As a check,
\[
\begin{pmatrix}a&b\\-b&a\end{pmatrix}
\begin{pmatrix}a&-b\\b&a\end{pmatrix}
=\begin{pmatrix}a^2+b^2&0\\0&a^2+b^2\end{pmatrix}
=(a^2+b^2)I_2,
\]
so, after division by $a^2+b^2\ne0$, the boxed matrix is the inverse
of $F(a+ib)$; if $a=b=0$, neither inverse exists.

\Needspace{5\baselineskip}
\par\bigskip
\noindent\textbf{2.2.4}\quad If $\mathsf{A}$ is an $n\times n$ matrix, show
that
\[
\det(-\mathsf{A})=(-1)^{n}\det\mathsf{A}.
\]
\solution
Changing $A$ into $-A$ multiplies each of its $n$ rows by $-1$.
Linearity of the determinant in each row gives $n$ such factors:
\[
\det(-A)=\underbrace{(-1)(-1)\cdots(-1)}_{n\ \text{rows}}\det A
=\boxed{(-1)^n\det A}.
\]
No nonsingularity assumption is needed.

\Needspace{5\baselineskip}
\par\bigskip
\noindent\textbf{2.2.5}\quad
\begin{enumerate}
\item[(a)] The matrix equation $\mathsf{A}^{2}=0$ does not imply
  $\mathsf{A}=0$. Show that the most general $2\times2$ matrix whose square is
  zero may be written as
  \[
  \begin{pmatrix}ab&b^{2}\\-a^{2}&-ab\end{pmatrix},
  \]
  where $a$ and $b$ are real or complex numbers.
  \erratum{Allow complex parameters for this complete parametrization,
  including when $\mathsf A$ is real. A parametrization using only real
  parameters needs the additional sign derived below.}
\item[(b)] If $\mathsf{C}=\mathsf{A}+\mathsf{B}$, in general
  \[
  \det\mathsf{C}\neq\det\mathsf{A}+\det\mathsf{B}.
  \]
  Construct a specific numerical example to illustrate this inequality.
\end{enumerate}
\solution
\textbf{(a)} Let $A=\begin{pmatrix}u&v\\w&t\end{pmatrix}$. Then
\[
A^2=\begin{pmatrix}u^2+vw&v(u+t)\\w(u+t)&t^2+vw\end{pmatrix}.
\]
If $v=w=0$, the equations $u^2=t^2=0$ force $A=0$,
represented by $a=b=0$. Otherwise one of $v,w$ is nonzero,
so the off-diagonal equations give $t=-u$. The remaining condition is
$u^2+vw=0$.
If $v\ne0$, choose a complex square root $b$ of $v$ and set
$a=u/b$. Then $ab=u$, $b^2=v$, and $-a^2=-u^2/v=w$.
If $v=0,w\ne0$, the condition gives $u=t=0$; choose $b=0$ and
$a^2=-w$. These cases prove completeness over $\mathbb C$.
Conversely,
\[
\begin{pmatrix}ab&b^2\\-a^2&-ab\end{pmatrix}^{\!2}
=\begin{pmatrix}
a^2b^2-a^2b^2&ab^3-ab^3\\
-a^3b+a^3b&-a^2b^2+a^2b^2
\end{pmatrix}=0.
\]
\erratum{For real matrices a complete real parametrization is
$s\begin{pmatrix}ab&b^2\\-a^2&-ab\end{pmatrix}$, with
$a,b\in\mathbb R$ and $s=\pm1$.}
For $v\ne0$, take $s=\operatorname{sgn}v$, $b=\sqrt{|v|}$,
$a=u/(sb)$; then $sab=u$, $sb^2=s|v|=v$, and
$-sa^2=-su^2/|v|=-u^2/v=w$ (using $s^2=1$ and $v=s|v|$).
For $v=0,w\ne0$, take
$s=-\operatorname{sgn}w$, $a=\sqrt{|w|}$, $b=0$; then
$-sa^2=\operatorname{sgn}(w)|w|=w$ and the other entries vanish.
For example $\begin{pmatrix}0&-1\\0&0\end{pmatrix}$ would require
$b^2=-1$ in the unmodified real formula.

\textbf{(b)} Take $A=B=I_2$. Then $\det A+\det B=2$ but
$\det(A+B)=\det(2I_2)=4$. Thus the determinant is not additive.

\Needspace{5\baselineskip}
\par\bigskip
\noindent\textbf{2.2.6}\quad Given
\[
\mathsf{K}=\begin{pmatrix}0&0&i\\-i&0&0\\0&-1&0\end{pmatrix},
\]
show that
\[
\mathsf{K}^{n}=\mathsf{KKK}\cdots(n\text{ factors})=\mathbf{1}
\]
(with the proper choice of $n$, $n\neq0$).
\solution
Multiplication gives
\[
K^2=\begin{pmatrix}0&-i&0\\0&0&1\\i&0&0\end{pmatrix},
\qquad K^3=\begin{pmatrix}-1&0&0\\0&-1&0\\0&0&-1\end{pmatrix}=-I_3.
\]
For example $(K^2)_{12}=K_{13}K_{32}=i(-1)=-i$; each row has
only one nonzero product. Multiplying $K^2$ by $K$ gives diagonal
entries $(-i)(-i)=-1$, $1(-1)=-1$, and $i(i)=-1$.
Therefore $\boxed{K^6=(K^3)^2=I_3}$.
This is the smallest positive exponent: $K,K^2,K^4=-K,K^5=-K^2$
are off-diagonal, and $K^3=-I_3\ne I_3$.
In general $K^n=I_3$ exactly when $n$ is a multiple of $6$.

\Needspace{5\baselineskip}
\par\bigskip
\noindent\textbf{2.2.7}\quad Verify the Jacobi identity,
\[
[\mathsf{A},[\mathsf{B},\mathsf{C}]]
=[\mathsf{B},[\mathsf{A},\mathsf{C}]]-[\mathsf{C},[\mathsf{A},\mathsf{B}]].
\]
\solution
Expand without reordering factors:
\[
\begin{aligned}
{}[A,[B,C]]&=ABC-ACB-BCA+CBA,\\
[B,[A,C]]&=BAC-BCA-ACB+CAB,\\
[C,[A,B]]&=CAB-CBA-ABC+BAC.
\end{aligned}
\]
Subtracting the third line from the second cancels $BAC$ and $CAB$:
\[
[B,[A,C]]-[C,[A,B]]=-BCA-ACB+CBA+ABC,
\]
which is the first line rearranged. This proves the identity for
arbitrary square matrices of the same size.

\Needspace{5\baselineskip}
\par\bigskip
\noindent\textbf{2.2.8}\quad Show that the matrices
\[
\mathsf{A}=\begin{pmatrix}0&1&0\\0&0&0\\0&0&0\end{pmatrix},\quad
\mathsf{B}=\begin{pmatrix}0&0&0\\0&0&1\\0&0&0\end{pmatrix},\quad
\mathsf{C}=\begin{pmatrix}0&0&1\\0&0&0\\0&0&0\end{pmatrix}
\]
satisfy the commutation relations
\[
[\mathsf{A},\mathsf{B}]=\mathsf{C},\quad[\mathsf{A},\mathsf{C}]=0,
\quad\text{and }[\mathsf{B},\mathsf{C}]=0.
\]
\solution
Let $E_{ij}$ have one $1$ at position $(i,j)$. Entrywise,
\[
(E_{ij}E_{kl})_{rs}
=\sum_t\delta_{ri}\delta_{tj}\delta_{tk}\delta_{sl}
=\delta_{jk}\delta_{ri}\delta_{sl},
\quad E_{ij}E_{kl}=\delta_{jk}E_{il}.
\]
Here $A=E_{12}$, $B=E_{23}$, $C=E_{13}$. Therefore
\[
\begin{array}{lll}
AB=E_{13}=C,&BA=0,&[A,B]=C,\\
AC=0,&CA=0,&[A,C]=0,\\
BC=0,&CB=0,&[B,C]=0.
\end{array}
\]
For example $BA=E_{23}E_{12}=0$ because $3\ne1$.

\Needspace{5\baselineskip}
\par\bigskip
\noindent\textbf{2.2.9}\quad Let
\[
\mathsf{i}=\begin{pmatrix}0&1&0&0\\-1&0&0&0\\0&0&0&1\\0&0&-1&0\end{pmatrix},\quad
\mathsf{j}=\begin{pmatrix}0&0&0&-1\\0&0&-1&0\\0&1&0&0\\1&0&0&0\end{pmatrix},
\]
and
\[
\mathsf{k}=\begin{pmatrix}0&0&-1&0\\0&0&0&1\\1&0&0&0\\0&-1&0&0\end{pmatrix}.
\]
Show that
\begin{enumerate}
\item[(a)] $\mathsf{i}^{2}=\mathsf{j}^{2}=\mathsf{k}^{2}=-\mathbf{1}$, where
  $\mathbf{1}$ is the unit matrix.
\item[(b)] $\mathsf{ij}=-\mathsf{ji}=\mathsf{k}$,\\
  $\mathsf{jk}=-\mathsf{kj}=\mathsf{i}$,\\
  $\mathsf{ki}=-\mathsf{ik}=\mathsf{j}$.
\end{enumerate}
These three matrices ($\mathsf{i}$, $\mathsf{j}$, and $\mathsf{k}$) plus the
unit matrix $\mathbf{1}$ form a basis for \textbf{quaternions}. An alternate
basis is provided by the four $2\times2$ matrices, $i\sigma_{1}$, $i\sigma_{2}$,
$-i\sigma_{3}$, and $\mathbf{1}$, where the $\sigma_{i}$ are the Pauli spin
matrices of Example 2.2.1.
\solution
\textbf{(a)} Track the four standard basis columns: column $r$ of a
matrix is its image of $e_r$. Reading off the columns, the given
matrices act as
\[
\begin{array}{c|rrrr}
&e_1&e_2&e_3&e_4\\\hline
\mathsf i&-e_2&e_1&-e_4&e_3\\
\mathsf j&e_4&e_3&-e_2&-e_1\\
\mathsf k&e_3&-e_4&-e_1&e_2
\end{array}
\]
Applying each row twice sends every $e_r$ to $-e_r$; hence
$\mathsf i^2=\mathsf j^2=\mathsf k^2=-I_4$.
For example $\mathsf i^2e_1=\mathsf i(-e_2)=-e_1$ and
$\mathsf j^2e_1=\mathsf j(e_4)=-e_1$.

\textbf{(b)} For the cyclic products apply the right-hand factor first,
e.g.\ $\mathsf i\mathsf je_1=\mathsf i(\mathsf je_1)=\mathsf ie_4=e_3$.
The successive images are
\[
\begin{array}{c|rrrr}
&e_1&e_2&e_3&e_4\\\hline
\mathsf i\mathsf j&e_3&-e_4&-e_1&e_2\\
\mathsf j\mathsf k&-e_2&e_1&-e_4&e_3\\
\mathsf k\mathsf i&e_4&e_3&-e_2&-e_1
\end{array}
\]
Thus these products equal $\mathsf k,\mathsf i,\mathsf j$.
For the reversed products, for example,
$\mathsf j\mathsf i=(\mathsf k\mathsf i)\mathsf i
=\mathsf k(-I_4)=-\mathsf k$.
Likewise $\mathsf k\mathsf j=(\mathsf i\mathsf j)\mathsf j=-\mathsf i$
and $\mathsf i\mathsf k=(\mathsf j\mathsf k)\mathsf k=-\mathsf j$.
The complete multiplication table is
\[
\begin{array}{c|rrr}
&\mathsf i&\mathsf j&\mathsf k\\\hline
\mathsf i&-I_4&\mathsf k&-\mathsf j\\
\mathsf j&-\mathsf k&-I_4&\mathsf i\\
\mathsf k&\mathsf j&-\mathsf i&-I_4
\end{array}.
\]

The matrices $I_4,\mathsf i,\mathsf j,\mathsf k$ are linearly
independent: the first row of
$a_0I_4+a_1\mathsf i+a_2\mathsf j+a_3\mathsf k$ is
$(a_0,a_1,-a_3,-a_2)$, which vanishes only if every $a_r=0$.
For the alternative basis, put $q_1=i\sigma_1$, $q_2=i\sigma_2$,
$q_3=-i\sigma_3$. The Pauli products of Exercise 2.2.11 give
$q_r^2=-\sigma_r^2=-I_2$ and
\[
q_1q_2=-\sigma_1\sigma_2=-i\sigma_3=q_3,\quad
q_2q_3=\sigma_2\sigma_3=i\sigma_1=q_1,\quad
q_3q_1=\sigma_3\sigma_1=i\sigma_2=q_2.
\]
The reversed products change sign. If
$a_0I_2+a_1q_1+a_2q_2+a_3q_3=0$, its entries are
\[
\begin{pmatrix}a_0-ia_3&ia_1+a_2\\ia_1-a_2&a_0+ia_3\end{pmatrix}=0.
\]
Adding and subtracting the diagonal equations gives $a_0=a_3=0$;
the off-diagonal equations give $a_1=a_2=0$. Thus these four
matrices are linearly independent and realize the same quaternion algebra
with real coefficients.

\Needspace{5\baselineskip}
\par\bigskip
\noindent\textbf{2.2.10}\quad A matrix with elements $a_{ij}=0$ for $j<i$ may
be called upper right triangular. The elements in the lower left (below and to
the left of the main diagonal) vanish. Show that the product of two upper
right triangular matrices is an upper right triangular matrix.
\solution
For $i>j$ the $(i,j)$ entry is $(AB)_{ij}=\sum_k a_{ik}b_{kj}$.
If $k<i$, then $a_{ik}=0$ by upper triangularity of $A$.
If $k\ge i$, then $k>j$ and $b_{kj}=0$ by upper triangularity
of $B$. Thus every summand is zero, and $(AB)_{ij}=0$ whenever
$i>j$. This is precisely upper triangularity of $AB$.

\Needspace{5\baselineskip}
\par\bigskip
\noindent\textbf{2.2.11}\quad The three Pauli spin matrices are
\[
\sigma_{1}=\begin{pmatrix}0&1\\1&0\end{pmatrix},\quad
\sigma_{2}=\begin{pmatrix}0&-i\\i&0\end{pmatrix},\quad\text{and}\quad
\sigma_{3}=\begin{pmatrix}1&0\\0&-1\end{pmatrix}.
\]
Show that
\begin{enumerate}
\item[(a)] $(\sigma_{i})^{2}=\mathbf{1}_{2}$,
\item[(b)] $\sigma_{i}\sigma_{j}=i\sigma_{k}$, $(i,j,k)=(1,2,3)$ or a cyclic
  permutation thereof,
\item[(c)] $\sigma_{i}\sigma_{j}+\sigma_{j}\sigma_{i}=2\delta_{ij}
  \mathbf{1}_{2}$; $\mathbf{1}_{2}$ is the $2\times2$ unit matrix.
\end{enumerate}
\solution
\textbf{(a)} Direct multiplication gives
\[
\sigma_1^2=I_2,\qquad
\sigma_2^2=\begin{pmatrix}(-i)i&0\\0&i(-i)\end{pmatrix}=I_2,\qquad
\sigma_3^2=\begin{pmatrix}1&0\\0&(-1)^2\end{pmatrix}=I_2.
\]

\textbf{(b)} The six products with distinct indices are
\[
\begin{array}{ll}
\sigma_1\sigma_2=\begin{pmatrix}i&0\\0&-i\end{pmatrix}=i\sigma_3,
&\sigma_2\sigma_1=-i\sigma_3,\\[3pt]
\sigma_2\sigma_3=\begin{pmatrix}0&i\\i&0\end{pmatrix}=i\sigma_1,
&\sigma_3\sigma_2=-i\sigma_1,\\[3pt]
\sigma_3\sigma_1=\begin{pmatrix}0&1\\-1&0\end{pmatrix}=i\sigma_2,
&\sigma_1\sigma_3=-i\sigma_2.
\end{array}
\]
The reversed products in the right column follow in the same way, or
from $\sigma_j\sigma_i=\sigma_j^\dagger\sigma_i^\dagger
=(\sigma_i\sigma_j)^\dagger=(i\sigma_k)^\dagger=-i\sigma_k$, since each
$\sigma_i$ is Hermitian and $(XY)^\dagger=Y^\dagger X^\dagger$
(Exercise 2.2.32); for example
$\sigma_2\sigma_1=\begin{pmatrix}-i&0\\0&i\end{pmatrix}=-i\sigma_3$.

\textbf{(c)} Hence the anticommutator is zero for distinct indices; for
equal indices it is $2\sigma_i^2=2I_2$. Together,
\[
\boxed{\sigma_i\sigma_j+\sigma_j\sigma_i=2\delta_{ij}I_2},
\qquad \sigma_i\sigma_j=\delta_{ij}I_2+
i\sum_k\varepsilon_{ijk}\sigma_k.
\]

\Needspace{5\baselineskip}
\par\bigskip
\noindent\textbf{2.2.12}\quad One description of spin-1 particles uses the
matrices
\[
\mathsf{M}_{x}=\frac{1}{\sqrt{2}}\begin{pmatrix}0&1&0\\1&0&1\\0&1&0
\end{pmatrix},\quad
\mathsf{M}_{y}=\frac{1}{\sqrt{2}}\begin{pmatrix}0&-i&0\\i&0&-i\\0&i&0
\end{pmatrix},
\]
and
\[
\mathsf{M}_{z}=\begin{pmatrix}1&0&0\\0&0&0\\0&0&-1\end{pmatrix}.
\]
Show that
\begin{enumerate}
\item[(a)] $[\mathsf{M}_{x},\mathsf{M}_{y}]=i\mathsf{M}_{z}$, and so on
  (cyclic permutation of indices). Using the Levi-Civita symbol, we may write
  \[
  [\mathsf{M}_{i},\mathsf{M}_{j}]=i\sum_{k}\varepsilon_{ijk}\mathsf{M}_{k}.
  \]
\item[(b)] $\mathsf{M}^{2}\equiv\mathsf{M}_{x}^{2}+\mathsf{M}_{y}^{2}
  +\mathsf{M}_{z}^{2}=2\,\mathbf{1}_{3}$, where $\mathbf{1}_{3}$ is the
  $3\times3$ unit matrix.
\item[(c)] $[\mathsf{M}^{2},\mathsf{M}_{i}]=0$,\\
  $[\mathsf{M}_{z},\mathsf{L}^{+}]=\mathsf{L}^{+}$,\\
  $[\mathsf{L}^{+},\mathsf{L}^{-}]=2\mathsf{M}_{z}$,\\
  where $\mathsf{L}^{+}\equiv\mathsf{M}_{x}+i\mathsf{M}_{y}$ and
  $\mathsf{L}^{-}\equiv\mathsf{M}_{x}-i\mathsf{M}_{y}$.
\end{enumerate}
\solution
\textbf{(a)} The two products involving $x,y$ are
\[
M_xM_y=\frac i2\begin{pmatrix}1&0&-1\\0&0&0\\1&0&-1\end{pmatrix},
\quad
M_yM_x=\frac i2\begin{pmatrix}-1&0&-1\\0&0&0\\1&0&1\end{pmatrix}.
\]
Their difference is $i\operatorname{diag}(1,0,-1)=iM_z$.
For the other cyclic pairs,
\[
\begin{aligned}
M_yM_z&=\frac i{\sqrt2}\begin{pmatrix}0&0&0\\1&0&1\\0&0&0\end{pmatrix},
&
M_zM_y&=\frac i{\sqrt2}\begin{pmatrix}0&-1&0\\0&0&0\\0&-1&0\end{pmatrix},\\
M_zM_x&=\frac1{\sqrt2}\begin{pmatrix}0&1&0\\0&0&0\\0&-1&0\end{pmatrix},
&
M_xM_z&=\frac1{\sqrt2}\begin{pmatrix}0&0&0\\1&0&-1\\0&0&0\end{pmatrix}.
\end{aligned}
\]
Subtraction gives
\[
[M_y,M_z]=\frac i{\sqrt2}\begin{pmatrix}0&1&0\\1&0&1\\0&1&0\end{pmatrix}
=iM_x,\qquad
[M_z,M_x]=\frac1{\sqrt2}\begin{pmatrix}0&1&0\\-1&0&1\\0&-1&0\end{pmatrix}
=iM_y,
\]
where the last matrix is $iM_y$ because $i(-i)=1$ and $i\cdot i=-1$.
Reversed commutators change sign and $[M_i,M_i]=0$, yielding the
Levi-Civita expression for all indices.

\textbf{(b)} The squares are
\[
M_x^2=\frac12\begin{pmatrix}1&0&1\\0&2&0\\1&0&1\end{pmatrix},
\quad
M_y^2=\frac12\begin{pmatrix}1&0&-1\\0&2&0\\-1&0&1\end{pmatrix},
\quad M_z^2=\operatorname{diag}(1,0,1).
\]
Adding cancels the off-diagonal entries $\pm\frac12$; the diagonal
entries add to $\frac12+\frac12+1=2$, $1+1+0=2$, and
$\frac12+\frac12+1=2$. Hence $\boxed{M^2=2I_3}$.

\textbf{(c)} Since $M^2=2I_3$, $[M^2,M_i]=2M_i-2M_i=0$.
Next, the matrix $iM_y$ found in (a) agrees with $M_x$ above the
diagonal and is its negative below it.
Adding and subtracting, the raising and lowering matrices are
\[
L^+=\sqrt2\begin{pmatrix}0&1&0\\0&0&1\\0&0&0\end{pmatrix},
\qquad L^-=\sqrt2\begin{pmatrix}0&0&0\\1&0&0\\0&1&0\end{pmatrix}.
\]
For a diagonal $D=\operatorname{diag}(d_1,d_2,d_3)$,
$[D,L]_{ij}=(d_i-d_j)L_{ij}$. On the nonzero entries of $L^+$
the difference for $M_z$ is $1$, so $[M_z,L^+]=L^+$.
Finally
\[
L^+L^-=\operatorname{diag}(2,2,0),\qquad
L^-L^+=\operatorname{diag}(0,2,2),
\]
whose difference is $\operatorname{diag}(2,0,-2)=2M_z$.

\Needspace{5\baselineskip}
\par\bigskip
\noindent\textbf{2.2.13}\quad Repeat Exercise 2.2.12, using the matrices for a
spin of $3/2$,
\[
\mathsf{M}_{x}=\frac{1}{2}\begin{pmatrix}
0&\sqrt{3}&0&0\\ \sqrt{3}&0&2&0\\ 0&2&0&\sqrt{3}\\ 0&0&\sqrt{3}&0
\end{pmatrix},\quad
\mathsf{M}_{y}=\frac{i}{2}\begin{pmatrix}
0&-\sqrt{3}&0&0\\ \sqrt{3}&0&-2&0\\ 0&2&0&-\sqrt{3}\\ 0&0&\sqrt{3}&0
\end{pmatrix},
\]
and
\[
\mathsf{M}_{z}=\frac{1}{2}\begin{pmatrix}
3&0&0&0\\ 0&1&0&0\\ 0&0&-1&0\\ 0&0&0&-3\end{pmatrix}.
\]
\erratum{In repeating part (b), determine the new value of $M^2$:
for these matrices the claim is $M^2=\tfrac{15}{4}I_4$.
The numerical spin-1 value is not retained.}
\solution
The commutator identities are unchanged, with $4\times4$ matrices.
All three parts use the ladder matrices $L^\pm=M_x\pm iM_y$, so form
them first. Since $i\cdot\frac i2=-\frac12$, the matrix $iM_y$ agrees
with $M_x$ above the diagonal and is its negative below it; hence
\[
L^+=M_x+iM_y=
\begin{pmatrix}0&\sqrt3&0&0\\0&0&2&0\\0&0&0&\sqrt3\\0&0&0&0\end{pmatrix},
\qquad L^-=M_x-iM_y=(L^+)^T.
\]
The diagonal entries of $M_z$ are $3/2,1/2,-1/2,-3/2$.
As in Exercise 2.2.12, $[M_z,L]_{ij}=((M_z)_{ii}-(M_z)_{jj})L_{ij}$
for any $L$. The nonzero entries of $L^+$ are at positions $(i,i+1)$,
where this difference is $1$, and those of $L^-$ are at $(i+1,i)$,
where it is $-1$. Thus
\[
[M_z,L^+]=L^+,\qquad [M_z,L^-]=-L^-.
\]
Because $L^-=(L^+)^T$ is real, $(L^+L^-)_{ii}$ is the sum of the
squares of row $i$ of $L^+$ and $(L^-L^+)_{ii}$ that of column $i$;
the off-diagonal entries vanish because each row and each column of
$L^+$ has at most one nonzero entry. Thus
\[
L^+L^-=\operatorname{diag}(3,4,3,0),\qquad
L^-L^+=\operatorname{diag}(0,3,4,3),
\]
so $[L^+,L^-]=\operatorname{diag}(3,1,-1,-3)=2M_z$.

\textbf{(a)} Now $M_x=(L^++L^-)/2$ and $M_y=(L^+-L^-)/(2i)$. Using
$[L^\pm,L^\pm]=0$, $[L^-,L^+]=-[L^+,L^-]$, and
$[M_z,L^\pm]=\pm L^\pm$,
\[
\begin{aligned}
{}[M_x,M_y]&=\frac1{4i}[L^++L^-,L^+-L^-]
=\frac1{4i}\bigl(-[L^+,L^-]+[L^-,L^+]\bigr)
=-\frac{4M_z}{4i}=iM_z,\\
[M_z,M_x]&=\tfrac12\bigl([M_z,L^+]+[M_z,L^-]\bigr)
=\tfrac12(L^+-L^-)=iM_y,\\
[M_y,M_z]&=-\tfrac1{2i}\bigl([M_z,L^+]-[M_z,L^-]\bigr)
=\tfrac1{2i}(-L^+-L^-)=iM_x.
\end{aligned}
\]
Reversed commutators change sign and $[M_i,M_i]=0$, so, as in
Exercise 2.2.12(a), $[M_i,M_j]=i\sum_k\varepsilon_{ijk}M_k$ for all
$i,j$.

\textbf{(b)} Expand the squares:
\[
\begin{aligned}
M_x^2+M_y^2
&=\tfrac14\bigl((L^+)^2+L^+L^-+L^-L^++(L^-)^2\bigr)\\
&\quad-\tfrac14\bigl((L^+)^2-L^+L^--L^-L^++(L^-)^2\bigr)\\
&=\tfrac12(L^+L^-+L^-L^+)
=\operatorname{diag}(3/2,7/2,7/2,3/2).
\end{aligned}
\]
Adding $M_z^2=\operatorname{diag}(9/4,1/4,1/4,9/4)$ gives, since
$\frac32+\frac94=\frac72+\frac14=\frac{15}4$,
$\boxed{M^2=\tfrac{15}{4}I_4}$.

\textbf{(c)} Since $M^2=\tfrac{15}{4}I_4$,
$[M^2,M_i]=\tfrac{15}{4}(M_i-M_i)=0$. The identities
$[M_z,L^+]=L^+$ and $[L^+,L^-]=2M_z$ were established at the start of
this solution.

\Needspace{5\baselineskip}
\par\bigskip
\noindent\textbf{2.2.14}\quad If $\mathsf{A}$ is a diagonal matrix, with all
diagonal elements different, and $\mathsf{A}$ and $\mathsf{B}$ commute, show
that $\mathsf{B}$ is diagonal.
\solution
Write $A=\operatorname{diag}(\lambda_1,\ldots,\lambda_n)$.
Then $(AB)_{ij}=\lambda_i b_{ij}$ and $(BA)_{ij}=b_{ij}\lambda_j$,
so $AB=BA$ implies
\[
(\lambda_i-\lambda_j)b_{ij}=0.
\]
For $i\ne j$ the hypothesis gives $\lambda_i-\lambda_j\ne0$,
so $b_{ij}=0$. Diagonal entries $b_{ii}$ are unrestricted.
Therefore $B$ is diagonal.

\Needspace{5\baselineskip}
\par\bigskip
\noindent\textbf{2.2.15}\quad If $\mathsf{A}$ and $\mathsf{B}$ are diagonal,
show that $\mathsf{A}$ and $\mathsf{B}$ commute.
\solution
For $A=\operatorname{diag}(a_1,\ldots,a_n)$ and
$B=\operatorname{diag}(b_1,\ldots,b_n)$,
\[
(AB)_{ij}=\sum_k a_i\delta_{ik}b_k\delta_{kj}=a_ib_i\delta_{ij}.
\]
Similarly $(BA)_{ij}=b_ia_i\delta_{ij}$. Scalar multiplication
commutes, so the entries agree and $AB=BA$.

\Needspace{5\baselineskip}
\par\bigskip
\noindent\textbf{2.2.16}\quad Show that
$\mathrm{trace}(\mathsf{ABC})=\mathrm{trace}(\mathsf{CBA})$ if any two of the
three matrices commute.
\solution
For square matrices of the same size,
\[
\operatorname{tr}(XY)=\sum_{i,j}x_{ij}y_{ji}
=\sum_{j,i}y_{ji}x_{ij}=\operatorname{tr}(YX).
\]
Taking $X=AB,Y=C$, and then $X=A,Y=BC$, gives the cyclic rule
$\operatorname{tr}(ABC)=\operatorname{tr}(BCA)=\operatorname{tr}(CAB)$.
There are three possible commuting pairs:
\[
\begin{array}{ll}
AB=BA:&\operatorname{tr}(ABC)=\operatorname{tr}(BAC)
=\operatorname{tr}(CBA),\\
BC=CB:&\operatorname{tr}(ABC)=\operatorname{tr}(ACB)
=\operatorname{tr}(CBA),\\
CA=AC:&\operatorname{tr}(ABC)=\operatorname{tr}(BCA)
=\operatorname{tr}(BAC)=\operatorname{tr}(CBA).
\end{array}
\]
Every possible hypothesis therefore gives the required equality.

\Needspace{5\baselineskip}
\par\bigskip
\noindent\textbf{2.2.17}\quad Angular momentum matrices satisfy a commutation
relation
\[
[\mathsf{M}_{j},\mathsf{M}_{k}]=i\mathsf{M}_{l},\quad j,k,l\text{ cyclic}.
\]
Show that the trace of each angular momentum matrix vanishes.
\erratum{Use finite matrices, so all products and traces below exist.}
\solution
For finite square matrices, cyclicity gives
$\operatorname{tr}[A,B]=\operatorname{tr}(AB)-\operatorname{tr}(BA)=0$.
Apply this to each cyclic angular momentum relation:
\[
\begin{aligned}
i\operatorname{tr}M_z&=\operatorname{tr}[M_x,M_y]=0,\\
i\operatorname{tr}M_x&=\operatorname{tr}[M_y,M_z]=0,\\
i\operatorname{tr}M_y&=\operatorname{tr}[M_z,M_x]=0.
\end{aligned}
\]
Division by $i$ proves that every trace vanishes. This uses ordinary
finite-dimensional traces, not undefined traces of infinite-dimensional
operators.

\Needspace{5\baselineskip}
\par\bigskip
\noindent\textbf{2.2.18}\quad $\mathsf{A}$ and $\mathsf{B}$ anticommute:
$\mathsf{AB}=-\mathsf{BA}$. Also, $\mathsf{A}^{2}=\mathbf{1}$,
$\mathsf{B}^{2}=\mathbf{1}$. Show that
$\mathrm{trace}(\mathsf{A})=\mathrm{trace}(\mathsf{B})=0$.

\emph{Note.}~The Pauli and Dirac matrices are specific examples.
\solution
From $AB=-BA$ and $B^2=I$, we obtain $BAB=-AB^2=-A$.
Taking traces and cycling the first $B$ to the end,
\[
-\operatorname{tr}A=\operatorname{tr}(BAB)
=\operatorname{tr}(ABB)=\operatorname{tr}A.
\]
Hence $2\operatorname{tr}A=0$.
Similarly $A^2=I$ gives $ABA=-BA^2=-B$, and cycling the first $A$
to the end,
$-\operatorname{tr}B=\operatorname{tr}(ABA)=\operatorname{tr}(BAA)
=\operatorname{tr}B$, so $2\operatorname{tr}B=0$.
Thus $\boxed{\operatorname{tr}A=\operatorname{tr}B=0}$.

\Needspace{5\baselineskip}
\par\bigskip
\noindent\textbf{2.2.19}\quad
\begin{enumerate}
\item[(a)] If two nonsingular matrices anticommute, show that the trace of
  each one is zero. (Nonsingular means that the determinant of the matrix is
  nonzero.)
\item[(b)] For the conditions of part (a) to hold, $\mathsf{A}$ and
  $\mathsf{B}$ must be $n\times n$ matrices with $n$ \textbf{even}. Show that
  if $n$ is \textbf{odd}, a contradiction results.
\end{enumerate}
\solution
\erratum{The Instructor's Manual's left multiplication by $B^{-1}$
proves the vanishing of $\operatorname{tr}A$, not of
$\operatorname{tr}B$ as printed.
Interchanging the roles of $A$ and $B$ supplies the second trace.}
\textbf{(a)} Multiplying $AB=-BA$ on the right by $B^{-1}$ gives
$A=-BAB^{-1}$, that is, $BAB^{-1}=-A$.
Taking the trace and cycling $B$ to the end,
\[
-\operatorname{tr}A=\operatorname{tr}(BAB^{-1})
=\operatorname{tr}(AB^{-1}B)=\operatorname{tr}A,
\]
so $\operatorname{tr}A=0$. Likewise, multiplying $BA=-AB$ on the
right by $A^{-1}$ gives $ABA^{-1}=-B$, and
$-\operatorname{tr}B=\operatorname{tr}(ABA^{-1})
=\operatorname{tr}(BA^{-1}A)=\operatorname{tr}B$.
Hence $\operatorname{tr}B=0$.

\textbf{(b)} Taking determinants and using the product rule and
$\det(-X)=(-1)^n\det X$ (Exercise 2.2.4) gives
\[
\det A\det B=\det(AB)=\det(-BA)
=(-1)^n\det B\det A.
\]
Both determinants are nonzero; division by their product yields
$1=(-1)^n$. Odd $n$ would give $1=-1$, a contradiction over the
real or complex numbers. Therefore $\boxed{n\text{ is even}}$.

\Needspace{5\baselineskip}
\par\bigskip
\noindent\textbf{2.2.20}\quad If $\mathsf{A}^{-1}$ has elements
\[
(\mathsf{A}^{-1})_{ij}=a_{ij}^{(-1)}=\frac{C_{ji}}{|\mathsf{A}|},
\]
where $C_{ji}$ is the $ji$th cofactor of $|\mathsf{A}|$, show that
\[
\mathsf{A}^{-1}\mathsf{A}=\mathbf{1}.
\]
Hence $\mathsf{A}^{-1}$ is the inverse of $\mathsf{A}$ (if
$|\mathsf{A}|\neq0$).
\solution
Let $B_{ij}=C_{ji}/\det A$, with $\det A\ne0$. Then
\[
(BA)_{ij}=\sum_k\frac{C_{ki}}{\det A}a_{kj}
=\frac1{\det A}\sum_k a_{kj}C_{ki}.
\]
For $i=j$, the numerator is the expansion down column $i$ and equals
$\det A$ (Exercise 2.1.4(a)). For $i\ne j$, it is the expansion of
the determinant with column $i$ replaced by column $j$, so it is zero
(Exercise 2.1.4(b)). Thus
$(BA)_{ij}=\delta_{ij}$ and $BA=I$.
The corresponding row identities of Exercise 2.1.4 give
\[
(AB)_{ij}=\frac1{\det A}\sum_k a_{ik}C_{jk}=\delta_{ij}.
\]
The formula therefore produces the two-sided inverse $B=A^{-1}$.

\Needspace{5\baselineskip}
\par\bigskip
\noindent\textbf{2.2.21}\quad Find the matrices $\mathsf{M}_{L}$ such that the
product $\mathsf{M}_{L}\mathsf{A}$ will be $\mathsf{A}$ but with:
\begin{enumerate}
\item[(a)] The $i$th row multiplied by a constant $k$ ($a_{ij}\to ka_{ij}$,
  $j=1,2,3,\ldots$);
\item[(b)] The $i$th row replaced by the original $i$th row minus a multiple
  of the $m$th row ($a_{ij}\to a_{ij}-Ka_{mj}$, $\erratum{j}=1,2,3,\ldots$);
\item[(c)] The $i$th and $m$th rows interchanged ($a_{ij}\to a_{mj}$,
  $a_{mj}\to a_{ij}$, $j=1,2,3,\ldots$).
\end{enumerate}
\solution
\erratum{In (b), the varying index in the parentheses must be $j$,
not $i$: $i$ labels the one row being changed.
For (c), the Instructor's Manual's $M_{mi}-M_{im}=1$ must read
$M_{mi}=M_{im}=1$ when $i\ne m$. Both inserted entries are positive.}
Let $E_{ab}$ have a single $1$ in position $(a,b)$. Its action is
\[
(E_{ab}A)_{rj}=\sum_s\delta_{ra}\delta_{sb}a_{sj}
=\delta_{ra}a_{bj}.
\]
Thus $E_{ab}A$ has row $b$ of $A$ in row $a$ and all other rows
zero. For distinct $i,m$, the requested matrices are
\[
\boxed{\begin{aligned}
\text{(a)}\quad M_L&=I+(k-1)E_{ii},\\
\text{(b)}\quad M_L&=I-KE_{im},\\
\text{(c)}\quad M_L&=I-E_{ii}-E_{mm}+E_{im}+E_{mi}.
\end{aligned}}
\]
In (a), row $i$ becomes $R_i+(k-1)R_i=kR_i$.
In (b), it becomes $R_i-KR_m$.
In (c), the two subtractions remove rows $i,m$; the last two
terms insert $R_m,R_i$ in their places. Every other row retains
only the identity contribution. If $i=m$, (b) reduces to scaling
by $1-K$, and the formula in (c) reduces to $I$.

\Needspace{5\baselineskip}
\par\bigskip
\noindent\textbf{2.2.22}\quad Find the matrices $\mathsf{M}_{R}$ such that the
product $\mathsf{AM}_{R}$ will be $\mathsf{A}$ but with:
\begin{enumerate}
\item[(a)] The $i$th column multiplied by a constant $k$ ($a_{ji}\to ka_{ji}$,
  $j=1,2,3,\ldots$);
\item[(b)] The $i$th column replaced by the original $i$th column minus a
  multiple of the $m$th column ($a_{ji}\to a_{ji}-ka_{jm}$, $j=1,2,3,\ldots$);
\item[(c)] The $i$th and $m$th columns interchanged ($a_{ji}\to a_{jm}$,
  $a_{jm}\to a_{ji}$, $j=1,2,3,\ldots$).
\end{enumerate}
\solution
For right multiplication,
\[
(AE_{ab})_{rj}=\sum_s a_{rs}\delta_{sa}\delta_{jb}
=a_{ra}\delta_{jb}.
\]
Thus $AE_{ab}$ puts column $a$ of $A$ in column $b$, giving
\[
\boxed{\begin{aligned}
\text{(a)}\quad M_R&=I+(k-1)E_{ii},\\
\text{(b)}\quad M_R&=I-kE_{mi},\\
\text{(c)}\quad M_R&=I-E_{ii}-E_{mm}+E_{im}+E_{mi}.
\end{aligned}}
\]
Column $i$ in (a) is $C_i+(k-1)C_i=kC_i$; in (b), it is
$C_i-kC_m$; in (c), the two columns are removed and reinserted
in exchanged positions.
\erratum{The off-diagonal entry in (b) is $(m,i)$, the transpose
of the row-operation entry $(i,m)$; the Instructor's Manual's
statement that the answers are the same requires this transpose.}
For rectangular $A$, $M_L$ has order equal to the number of rows
of $A$, whereas $M_R$ has order equal to the number of columns.

\Needspace{5\baselineskip}
\par\bigskip
\noindent\textbf{2.2.23}\quad Find the inverse of
\[
\mathsf{A}=\begin{pmatrix}3&2&1\\2&2&1\\1&1&4\end{pmatrix}.
\]
\solution
Expansion along row 1 gives
\[
\det A=3(8-1)-2(8-1)+(2-2)=21-14=7.
\]
The cofactors (signed $2\times2$ minors) are
\[
\begin{array}{lll}
C_{11}=8-1=7,&C_{12}=-(8-1)=-7,&C_{13}=2-2=0,\\
C_{21}=-(8-1)=-7,&C_{22}=12-1=11,&C_{23}=-(3-2)=-1,\\
C_{31}=2-2=0,&C_{32}=-(3-2)=-1,&C_{33}=6-4=2.
\end{array}
\]
By Exercise 2.2.20, $(A^{-1})_{ij}=C_{ji}/\det A$, so $A^{-1}$ is the
transposed cofactor matrix divided by $7$. The cofactor matrix is
symmetric, so its transpose is unchanged:
\[
\boxed{A^{-1}=\frac17
\begin{pmatrix}7&-7&0\\-7&11&-1\\0&-1&2\end{pmatrix}}.
\]
An independent multiplication checks every entry:
\[
A\begin{pmatrix}7&-7&0\\-7&11&-1\\0&-1&2\end{pmatrix}
=\begin{pmatrix}
21-14&-21+22-1&-2+2\\
14-14&-14+22-1&-2+2\\
7-7&-7+11-4&-1+8
\end{pmatrix}=7I_3.
\]

\Needspace{5\baselineskip}
\par\bigskip
\noindent\textbf{2.2.24}\quad Matrices are far too useful to remain the
exclusive property of physicists. They may appear wherever there are linear
relations. For instance, in a study of population movement the initial
fraction of a fixed population in each of $n$ areas (or industries or
religions, etc.) is represented by an $n$-component column vector
$\mathsf{P}$. The movement of people from one area to another in a given time
is described by an $n\times n$ (stochastic) matrix $\mathsf{T}$. Here
$T_{ij}$ is the fraction of the population in the $j$th area that moves to the
$i$th area. (Those not moving are covered by $i=j$.) With $\mathsf{P}$
describing the initial population distribution, the final population
distribution is given by the matrix equation $\mathsf{TP}=\mathsf{Q}$. From
its definition, $\sum_{i=1}^{n}P_{i}=1$.
\begin{enumerate}
\item[(a)] Show that conservation of people requires that
  \[
  \sum_{i=1}^{n}T_{ij}=1,\quad j=1,2,\ldots,n.
  \]
\item[(b)] Prove that
  \[
  \sum_{i=1}^{n}Q_{i}=1
  \]
  continues the conservation of people.
\end{enumerate}
\solution
\textbf{(a)} A source area $j$ containing $N_j$ people contributes
$T_{1j}N_j,\ldots,T_{nj}N_j$ to the destination areas.
Conservation requires the sum to be $N_j$; for $N_j>0$, division
gives $\sum_iT_{ij}=1$. This requirement holds for every source
area, even one initially empty, because conservation must hold for
every possible initial distribution. Equivalently, start with $P=e_j$;
the final total $\sum_iT_{ij}$ must be $1$.

\textbf{(b)} Since $Q_i=\sum_jT_{ij}P_j$,
\[
\sum_iQ_i=\sum_i\sum_jT_{ij}P_j
=\sum_jP_j\left(\sum_iT_{ij}\right)=\sum_jP_j=1.
\]
Nonnegative fractions and $P_j\ge0$ also imply $Q_i\ge0$, so $Q$
is again a valid population distribution.

\Needspace{5\baselineskip}
\par\bigskip
\noindent\textbf{2.2.25}\quad Given a $6\times6$ matrix $\mathsf{A}$ with
elements $a_{ij}=0.5^{|i-j|}$, $i,j=0,1,2,\ldots,5$, find $\mathsf{A}^{-1}$.
\par\nobreak\hfill\emph{ANS.}\enspace
$\mathsf{A}^{-1}=\dfrac{1}{3}\begin{pmatrix}
4&-2&0&0&0&0\\-2&5&-2&0&0&0\\0&-2&5&-2&0&0\\
0&0&-2&5&-2&0\\0&0&0&-2&5&-2\\0&0&0&0&-2&4
\end{pmatrix}$.
\solution
First retain a parameter $r$ in $A_{ij}=r^{|i-j|}$ for
$i,j=0,\ldots,5$. Subtract $r$ times original row $i-1$ from
row $i$ for each $i=1,\ldots,5$. For $j<i$, the new entry is
$r^{i-j}-r\,r^{i-1-j}=0$; for $j\ge i$, it is
$r^{j-i}-r\,r^{j-i+1}=(1-r^2)r^{j-i}$.
The result is $EA$, where $E$ has $1$ on the diagonal and $-r$ just
below it. Now subtract $r$ times original column $j-1$ of $EA$ from
column $j$, for each $j=1,\ldots,5$, which is right multiplication by
$E^T$. Above the
diagonal this gives $r^j-r\,r^{j-1}=0$ in row $0$ and
$(1-r^2)(r^{j-i}-r\,r^{j-1-i})=0$ for $j>i\ge1$; the diagonal and
lower entries are unchanged, because the entry to the left of each is
$0$. Hence
\[
EAE^T=\Delta\equiv\operatorname{diag}(1,1-r^2,\ldots,1-r^2).
\]
Since $E$ is lower triangular with unit diagonal, $\det E=1$, so
$\det A=\det\Delta=(1-r^2)^5$: $A$ is invertible exactly when
$r^2\ne1$ (for $r=\pm1$ every row of $A$ is $\pm$ row~$0$).
For $r^2\ne1$, $A=E^{-1}\Delta(E^T)^{-1}$ gives
$A^{-1}=E^T\Delta^{-1}E$.
Write $\Delta^{-1}=\operatorname{diag}(\delta_0,\ldots,\delta_5)$, so
$\delta_0=1$ and $\delta_1=\cdots=\delta_5=1/(1-r^2)$, and set
$\delta_6=0$. Column $i$ of $E$ has only the entries $1$ (row $i$) and
$-r$ (row $i+1$), so $(E^T\Delta^{-1}E)_{ij}=\sum_kE_{ki}\delta_kE_{kj}$
receives contributions only from $k\in\{i,i+1\}\cap\{j,j+1\}$. It is
$\delta_i+r^2\delta_{i+1}$ on the diagonal, $-r\delta_{\max(i,j)}$ for
$|i-j|=1$, and zero elsewhere. Explicitly,
\[
\begin{aligned}
\delta_0+r^2\delta_1&=1+\frac{r^2}{1-r^2}=\frac1{1-r^2},&
\delta_5+r^2\delta_6&=\frac1{1-r^2},\\
\delta_i+r^2\delta_{i+1}&=\frac{1+r^2}{1-r^2}\quad(1\le i\le4),&
-r\delta_{\max(i,j)}&=-\frac{r}{1-r^2}.
\end{aligned}
\]
This is $T/(1-r^2)$, where $T$ is
tridiagonal with end diagonal entries $1$, interior diagonal entries
$1+r^2$, and adjacent entries $-r$.
The result can be confirmed by multiplying $A$ and $T$ directly.
For an interior column $j=1,\ldots,4$,
\[
(AT)_{ij}=-rA_{i,j-1}+(1+r^2)A_{ij}-rA_{i,j+1}.
\]
If $i<j$, put $d=j-i\ge1$ to obtain
$-r r^{d-1}+(1+r^2)r^d-r r^{d+1}=0$.
If $i>j$, put $d=i-j\ge1$ to obtain
$-r r^{d+1}+(1+r^2)r^d-r r^{d-1}=0$.
For $i=j$, the result is $-r^2+(1+r^2)-r^2=1-r^2$.
At the first column,
\[
(AT)_{i0}=A_{i0}-rA_{i1}
=\begin{cases}1-r^2,&i=0,\\r^i-r r^{i-1}=0,&i>0.\end{cases}
\]
At the last column, likewise,
\[
(AT)_{i5}=-rA_{i4}+A_{i5}
=\begin{cases}1-r^2,&i=5,\\-r\,r^{4-i}+r^{5-i}=0,&i<5.\end{cases}
\]
Thus $AT=(1-r^2)I_6$, confirming $A^{-1}=T/(1-r^2)$.
At $r=1/2$, $1-r^2=3/4$, so $A^{-1}=\frac43T$ has end diagonal
entries $\frac43$, interior diagonal entries $\frac43\cdot\frac54=\frac53$,
and adjacent entries $\frac43\bigl(-\frac12\bigr)=-\frac23$:
\[
\boxed{A^{-1}=\frac{T}{1-r^2}
=\frac13\begin{pmatrix}
4&-2&0&0&0&0\\
-2&5&-2&0&0&0\\
0&-2&5&-2&0&0\\
0&0&-2&5&-2&0\\
0&0&0&-2&5&-2\\
0&0&0&0&-2&4
\end{pmatrix}}.
\]

\Needspace{5\baselineskip}
\par\bigskip
\noindent\textbf{2.2.26}\quad Show that the product of two orthogonal matrices
is orthogonal.
\solution
Let $A^TA=I$ and $B^TB=I$. Since transposition reverses products,
\[
(AB)^T(AB)=B^TA^TAB=B^T(A^TA)B=B^TIB=B^TB=I.
\]
This is exactly the orthogonality condition for $AB$.

\Needspace{5\baselineskip}
\par\bigskip
\noindent\textbf{2.2.27}\quad If $\mathsf{A}$ is orthogonal, show that its
determinant $=\pm1$.
\solution
Taking determinants of $A^TA=I$ and using $\det(A^T)=\det A$ gives
\[
1=\det(A^TA)=\det(A^T)\det A=(\det A)^2.
\]
The two roots of $d^2=1$ are $1,-1$, so $\boxed{\det A=\pm1}$.

\Needspace{5\baselineskip}
\par\bigskip
\noindent\textbf{2.2.28}\quad Show that the trace of the product of a
symmetric and an antisymmetric matrix is zero.
\solution
Let $S^T=S$ and $K^T=-K$. Transposition preserves the trace:
\[
\operatorname{tr}(SK)=\operatorname{tr}((SK)^T)
=\operatorname{tr}(K^TS^T)
=-\operatorname{tr}(KS)=-\operatorname{tr}(SK).
\]
Thus $2\operatorname{tr}(SK)=0$ and $\boxed{\operatorname{tr}(SK)=0}$.
Entrywise, each pair with $i<j$ cancels:
$s_{ij}k_{ji}+s_{ji}k_{ij}=s_{ij}(-k_{ij})+s_{ij}k_{ij}=0$,
while all diagonal terms vanish because $k_{ii}=0$.

\Needspace{5\baselineskip}
\par\bigskip
\noindent\textbf{2.2.29}\quad $\mathsf{A}$ is $2\times2$ and orthogonal. Find
the most general form of
\[
\mathsf{A}=\begin{pmatrix}a&b\\c&d\end{pmatrix}.
\]
\solution
For real orthogonal $A$, $A^TA=I$ means $A^T=A^{-1}$, so also
$AA^T=I$: the rows are orthonormal. The first row is a real unit
vector, so write $(a,b)=(\cos\theta,\sin\theta)$.
The second row is a unit vector perpendicular to the first.
The vectors perpendicular to $(\cos\theta,\sin\theta)$ form the line
spanned by $(-\sin\theta,\cos\theta)$, whose only unit vectors are
$\pm(-\sin\theta,\cos\theta)$. Hence
\[
\boxed{A=\begin{pmatrix}
\cos\theta&\sin\theta\\
-s\sin\theta&s\cos\theta
\end{pmatrix},\quad s=\pm1,\quad\theta\in\mathbb R}.
\]
The two row norms are $1$, their inner product is $0$, and
$\det A=s(\cos^2\theta+\sin^2\theta)=s$.
For $s=1$, $A$ is the two-dimensional rotation matrix: it rotates the
coordinate axes through the angle $\theta$,
$x'=x\cos\theta+y\sin\theta$, $y'=-x\sin\theta+y\cos\theta$.
For $s=-1$, $A=\operatorname{diag}(1,-1)$ times that rotation, i.e.\
the rotation followed by the reflection $y'\to-y'$.
\erratum{Both signs are required. The Instructor's Manual assumes
$\det A=1$ and omits the $\det A=-1$ reflection family.}

If orthogonality is extended to complex $A$ with $A^TA=I$, the
algebraic form is
$A=\begin{pmatrix}a&b\\-sb&sa\end{pmatrix}$ with
$a^2+b^2=1$, $a,b\in\mathbb C$, and $s=\pm1$.
Indeed, as in the real case $A^TA=I$ gives $AA^T=I$, so
$a^2+b^2=1$, $ac+bd=0$, and $c^2+d^2=1$. Since $(a,b)\ne(0,0)$, the
solutions $(c,d)$ of the single linear equation $ac+bd=0$ form a
one-dimensional space containing $(-b,a)\ne0$, so
$(c,d)=\lambda(-b,a)$; then $c^2+d^2=\lambda^2(a^2+b^2)=\lambda^2=1$
gives $\lambda=s=\pm1$. Conversely, every such $A$ has
$A^TA=(a^2+b^2)I=I$.

\Needspace{5\baselineskip}
\par\bigskip
\noindent\textbf{2.2.30}\quad Show that
\[
\det(\mathsf{A}^{*})=(\det\mathsf{A})^{*}=\det(\mathsf{A}^{\dagger}).
\]
\solution
The permutation formula is
\[
\det A=\sum_{\pi\in S_n}\operatorname{sgn}(\pi)
\prod_{r=1}^n a_{r,\pi(r)}.
\]
Because each permutation sign is real, conjugating the finite sum gives
\[
(\det A)^*=\sum_\pi\operatorname{sgn}(\pi)
\prod_r a_{r,\pi(r)}^*=\det(A^*).
\]
Finally $A^\dagger=(A^*)^T$, and determinants are invariant under
transposition. Thus $\det(A^\dagger)=\det(A^*)=(\det A)^*$.

\Needspace{5\baselineskip}
\par\bigskip
\noindent\textbf{2.2.31}\quad Three angular momentum matrices satisfy the
basic commutation relation
\[
[\mathsf{J}_{x},\mathsf{J}_{y}]=i\mathsf{J}_{z}
\]
(and cyclic permutation of indices). If two of the matrices have real
elements, show that the elements of the third must be pure imaginary.
\solution
Suppose the real matrices are $J_x,J_y$. Each entry
\[
R_{ab}:=([J_x,J_y])_{ab}
=\sum_c\bigl((J_x)_{ac}(J_y)_{cb}-(J_y)_{ac}(J_x)_{cb}\bigr)
\]
is real. The commutation relation gives
$(J_z)_{ab}=R_{ab}/i=-iR_{ab}$, which is pure imaginary
(possibly zero). For either other real pair, use the corresponding
cyclic commutation relation in the same calculation.

\Needspace{5\baselineskip}
\par\bigskip
\noindent\textbf{2.2.32}\quad Show that
$(\mathsf{AB})^{\dagger}=\mathsf{B}^{\dagger}\mathsf{A}^{\dagger}$.
\solution
For conformable matrices,
\[
\begin{aligned}
((AB)^\dagger)_{ij}
&=(AB)_{ji}^*
=\left(\sum_k a_{jk}b_{ki}\right)^*
=\sum_k a_{jk}^*b_{ki}^*\\
&=\sum_k(B^\dagger)_{ik}(A^\dagger)_{kj}
=(B^\dagger A^\dagger)_{ij}.
\end{aligned}
\]
The scalar entries commute; no commutation of the matrices was
assumed. Equality of all entries proves the reversed product rule.

\Needspace{5\baselineskip}
\par\bigskip
\noindent\textbf{2.2.33}\quad A matrix
$\mathsf{C}=\mathsf{S}^{\dagger}\mathsf{S}$. Show that the trace is positive
definite unless $\mathsf{S}$ is the null matrix, in which case
$\mathrm{trace}\,(\mathsf{C})=0$.
\solution
If $S$ has entries $s_{ki}$, the diagonal entries of $C=S^\dagger S$ are
\[
C_{ii}=\sum_k(S^\dagger)_{ik}S_{ki}
=\sum_k s_{ki}^*s_{ki}=\sum_k|s_{ki}|^2.
\]
Therefore
\[
\boxed{\operatorname{tr}(S^\dagger S)=\sum_{i,k}|s_{ki}|^2}.
\]
Every summand is nonnegative. The sum is zero exactly when every
entry is zero; otherwise it is strictly positive. The trace is thus
a positive-definite quadratic form in the real and imaginary parts of
the entries of $S$.
This does not imply that $C$ itself is positive definite:
$S=\operatorname{diag}(1,0)$ gives $C=\operatorname{diag}(1,0)$,
which has a zero eigenvalue.

\Needspace{5\baselineskip}
\par\bigskip
\noindent\textbf{2.2.34}\quad If $\mathsf{A}$ and $\mathsf{B}$ are Hermitian
matrices, show that $(\mathsf{AB}+\mathsf{BA})$ and $i(\mathsf{AB}-\mathsf{BA})$
are also Hermitian.
\solution
Using $A^\dagger=A$, $B^\dagger=B$, and reversing the product order
(Exercise 2.2.32),
\[
(AB+BA)^\dagger=B^\dagger A^\dagger+A^\dagger B^\dagger
=BA+AB=AB+BA.
\]
For the second expression also conjugate the scalar $i$:
\[
\bigl(i(AB-BA)\bigr)^\dagger
=-i(B^\dagger A^\dagger-A^\dagger B^\dagger)
=-i(BA-AB)=i(AB-BA).
\]
Both expressions equal their adjoints and are therefore Hermitian.

\Needspace{5\baselineskip}
\par\bigskip
\noindent\textbf{2.2.35}\quad The matrix $\mathsf{C}$ is not Hermitian. Show
that then $\mathsf{C}+\mathsf{C}^{\dagger}$ and
$i(\mathsf{C}-\mathsf{C}^{\dagger})$ are Hermitian. This means that a
non-Hermitian matrix may be resolved into two Hermitian parts,
\[
\mathsf{C}=\frac{1}{2}(\mathsf{C}+\mathsf{C}^{\dagger})
+\frac{1}{2i}\,i(\mathsf{C}-\mathsf{C}^{\dagger}).
\]
This decomposition of a matrix into two Hermitian matrix parts parallels the
decomposition of a complex number $z$ into $x+iy$, where $x=(z+z^{*})/2$ and
$y=(z-z^{*})/2i$.
\solution
For any square $C$, whether or not it is Hermitian, set
\[
H=\frac{C+C^\dagger}{2},\qquad K=\frac{C-C^\dagger}{2i}.
\]
Taking adjoints with $(C^\dagger)^\dagger=C$ and
$(\lambda X)^\dagger=\lambda^*X^\dagger$, where $(1/2i)^*=1/(-2i)$,
\[
H^\dagger=\frac{C^\dagger+C}{2}=H,\qquad
K^\dagger=\frac{C^\dagger-C}{-2i}
=\frac{C-C^\dagger}{2i}=K.
\]
Also $C+C^\dagger=2H$ and $i(C-C^\dagger)=i\cdot2iK=-2K$. Real
multiples of Hermitian matrices are Hermitian, so both matrices
requested in the first sentence are Hermitian.
Recombining, with $iK=i(C-C^\dagger)/(2i)=(C-C^\dagger)/2$,
\[
H+iK=\frac{C+C^\dagger}{2}+\frac{C-C^\dagger}{2}=C.
\]
This is the displayed decomposition: its second summand is the
Hermitian matrix $i(C-C^\dagger)=-2K$ times the scalar $1/(2i)$, that
is, $iK$. The summand $iK=(C-C^\dagger)/2$ itself is anti-Hermitian,
just as $iy$ is imaginary; the two Hermitian parts are $H$ and $K$,
the matrix analogues of $x=\operatorname{Re}z$ and
$y=\operatorname{Im}z$.
The decomposition is unique: if $C=H_1+iK_1$ with Hermitian
$H_1,K_1$, then $C^\dagger=H_1-iK_1$; adding and subtracting give
$C+C^\dagger=2H_1$ and $C-C^\dagger=2iK_1$, so $H_1=H$ and $K_1=K$.

\Needspace{5\baselineskip}
\par\bigskip
\noindent\textbf{2.2.36}\quad $\mathsf{A}$ and $\mathsf{B}$ are two
noncommuting Hermitian matrices:
\[
\mathsf{AB}-\mathsf{BA}=i\mathsf{C}.
\]
Prove that $\mathsf{C}$ is Hermitian.
\solution
Take adjoints of both sides of $AB-BA=iC$. By the reversed product
rule $(XY)^\dagger=Y^\dagger X^\dagger$ (Exercise 2.2.32) and
$A^\dagger=A$, $B^\dagger=B$, the left side gives
\[
(AB-BA)^\dagger=B^\dagger A^\dagger-A^\dagger B^\dagger
=BA-AB=-(AB-BA)=-iC,
\]
whereas the right side gives $(iC)^\dagger=i^*C^\dagger=-iC^\dagger$.
Therefore $-iC^\dagger=-iC$, and division by $-i$ proves
$\boxed{C^\dagger=C}$.

\Needspace{5\baselineskip}
\par\bigskip
\noindent\textbf{2.2.37}\quad Two matrices $\mathsf{A}$ and $\mathsf{B}$ are
each Hermitian. Find a necessary and sufficient condition for their product
$\mathsf{AB}$ to be Hermitian.
\par\nobreak\hfill\emph{ANS.}\enspace $[\mathsf{A},\mathsf{B}]=0$.
\solution
By the reversed product rule (Exercise 2.2.32) and Hermiticity of
$A,B$, $(AB)^\dagger=B^\dagger A^\dagger=BA$.
If $AB$ is Hermitian, then $AB=(AB)^\dagger=BA$; conversely, if
$AB=BA$, then $(AB)^\dagger=BA=AB$. Thus $AB$ is Hermitian precisely
when $AB=BA$, or $\boxed{[A,B]=0}$. This proves necessity and
sufficiency.

\Needspace{5\baselineskip}
\par\bigskip
\noindent\textbf{2.2.38}\quad Show that the reciprocal (that is, inverse) of a
unitary matrix is unitary.
\solution
Unitary $U$ satisfies $U^{-1}=U^\dagger$. Put $V=U^{-1}$.
Then $V^\dagger=(U^\dagger)^\dagger=U$, so
\[
V^\dagger V=UU^{-1}=I,\qquad VV^\dagger=U^{-1}U=I.
\]
Hence $V$ is unitary.

\Needspace{5\baselineskip}
\par\bigskip
\noindent\textbf{2.2.39}\quad Prove that the direct product of two unitary
matrices is unitary.
\solution
\erratum{``Direct product'' here means the tensor product defined
in the text. The Instructor's Manual instead displays the proof for
the ordinary product $UV$; the double-index proof below supplies the
requested result.}
Let unitary $U,V$ have orders $m,n$. Their direct (Kronecker)
product $W=U\otimes V$ has rows and columns labeled by index pairs,
with entries $W_{(i,a),(j,b)}=U_{ij}V_{ab}$. Its adjoint has entries
$(W^\dagger)_{(j,b),(i,a)}=W^*_{(i,a),(j,b)}=U_{ij}^*V_{ab}^*$, and a
matrix product sums over the intermediate pair $(i,a)$. Thus
\[
\begin{aligned}
(W^\dagger W)_{(j,b),(k,c)}
&=\sum_{i=1}^m\sum_{a=1}^n U_{ij}^*V_{ab}^*U_{ik}V_{ac}
=\left(\sum_iU_{ij}^*U_{ik}\right)
\left(\sum_aV_{ab}^*V_{ac}\right)\\
&=(U^\dagger U)_{jk}(V^\dagger V)_{bc}=\delta_{jk}\delta_{bc}.
\end{aligned}
\]
Therefore $W^\dagger W=I_{mn}$. In the reverse order the sum runs
over the pair $(j,b)$:
\[
(WW^\dagger)_{(i,a),(k,c)}
=\sum_{j,b}U_{ij}V_{ab}U_{kj}^*V_{cb}^*
=(UU^\dagger)_{ik}(VV^\dagger)_{ac}=\delta_{ik}\delta_{ac},
\]
so $WW^\dagger=I_{mn}$ as well, proving unitarity of the direct
product. In matrix notation, the entry formula for $W^\dagger$ says
$(U\otimes V)^\dagger=U^\dagger\otimes V^\dagger$, and with the
direct-product rule $(A\otimes B)(C\otimes D)=AC\otimes BD$ the same
calculation reads
\[
(U\otimes V)^\dagger(U\otimes V)
=(U^\dagger U)\otimes(V^\dagger V)=I_m\otimes I_n=I_{mn}.
\]

\Needspace{5\baselineskip}
\par\bigskip
\noindent\textbf{2.2.40}\quad If $\boldsymbol{\sigma}$ is the vector with the
$\sigma_{i}$ as components given in Eq.~(2.61), and $\mathbf{p}$ is an ordinary
vector, show that
\[
(\boldsymbol{\sigma}\cdot\mathbf{p})^{2}=\mathbf{p}^{2}\mathbf{1}_{2},
\]
where $\mathbf{1}_{2}$ is a $2\times2$ unit matrix.
\solution
\erratum{In the Instructor's Manual's expansion, the term multiplying
$p_yp_z$ must contain $\sigma_2\sigma_3+\sigma_3\sigma_2$, not the
repeated $\sigma_1\sigma_2+\sigma_2\sigma_1$. Both vanish, so the
printed final result is unaffected.}
The scalar components $p_i$ commute with each other and with the
matrices, so
\[
(\boldsymbol{\sigma}\cdot\mathbf p)^2
=\Bigl(\sum_ip_i\sigma_i\Bigr)\Bigl(\sum_jp_j\sigma_j\Bigr)
=\sum_{i,j}p_ip_j\,\sigma_i\sigma_j.
\]
Collect the three terms with $i=j$ and, for each pair $i<j$, the
two terms $(i,j)$ and $(j,i)$:
\[
\begin{aligned}
(\boldsymbol{\sigma}\cdot\mathbf p)^2
={}&p_1^2\sigma_1^2+p_2^2\sigma_2^2+p_3^2\sigma_3^2\\
&+p_1p_2(\sigma_1\sigma_2+\sigma_2\sigma_1)
+p_1p_3(\sigma_1\sigma_3+\sigma_3\sigma_1)\\
&+p_2p_3(\sigma_2\sigma_3+\sigma_3\sigma_2).
\end{aligned}
\]
By Exercise 2.2.11, each square is $\sigma_i^2=I_2$ and each mixed
anticommutator $\sigma_i\sigma_j+\sigma_j\sigma_i$ ($i\ne j$) is zero:
\[
\boxed{(\boldsymbol{\sigma}\cdot\mathbf p)^2
=(p_1^2+p_2^2+p_3^2)I_2}.
\]
The same reasoning applies to the mutually commuting free-particle
momentum operators $-i\hbar\partial_i$ on functions for which the mixed
derivatives commute.

\Needspace{5\baselineskip}
\par\bigskip
\noindent\textbf{2.2.41}\quad Use the equations for the properties of direct
products, Eqs.~(2.57) and (2.58), to show that the four matrices
$\boldsymbol{\gamma}^{\mu}$, $\mu=0,1,2,3$, satisfy the conditions listed in
Eqs.~(2.74) and (2.75).
\solution
Explicitly, the Dirac matrices of Eqs.~(2.70) and (2.72) are
\[
\gamma^0=\sigma_3\otimes I_2,\qquad
\gamma^i=J\otimes\sigma_i\quad(i=1,2,3),\qquad
J=\begin{pmatrix}0&1\\-1&0\end{pmatrix}=i\sigma_2.
\]
The direct-product rule is $(A\otimes B)(C\otimes D)=AC\otimes BD$.
Since $J^2=(i\sigma_2)^2=-\sigma_2^2=-I_2$ and $\sigma_i^2=I_2$
(Exercise 2.2.11),
\[
(\gamma^0)^2=\sigma_3^2\otimes I_2^2=I_2\otimes I_2=I_4,\qquad
(\gamma^i)^2=J^2\otimes\sigma_i^2=-I_2\otimes I_2=-I_4.
\]
For the mixed products the rule gives
$\gamma^0\gamma^i=(\sigma_3\otimes I_2)(J\otimes\sigma_i)
=\sigma_3J\otimes\sigma_i$ and
$\gamma^i\gamma^0=J\sigma_3\otimes\sigma_i$, where
\[
\sigma_3J=\begin{pmatrix}1&0\\0&-1\end{pmatrix}
\begin{pmatrix}0&1\\-1&0\end{pmatrix}
=\begin{pmatrix}0&1\\1&0\end{pmatrix}=\sigma_1,\qquad
J\sigma_3=\begin{pmatrix}0&-1\\-1&0\end{pmatrix}=-\sigma_1.
\]
Hence
\[
\gamma^0\gamma^i+\gamma^i\gamma^0
=(\sigma_3J+J\sigma_3)\otimes\sigma_i=0.
\]
For spatial indices $\gamma^i\gamma^j=J^2\otimes\sigma_i\sigma_j
=-I_2\otimes\sigma_i\sigma_j$, so by Exercise 2.2.11(c)
\[
\gamma^i\gamma^j+\gamma^j\gamma^i
=-I_2\otimes(\sigma_i\sigma_j+\sigma_j\sigma_i)
=-I_2\otimes2\delta_{ij}I_2=-2\delta_{ij}I_4.
\]
For $i=j$ this repeats $(\gamma^i)^2=-I_4$; for $i\ne j$ it is the
anticommutation of two spatial matrices.
These are the conditions (2.74)--(2.75), compactly
\[
\boxed{\{\gamma^\mu,\gamma^\nu\}=2g^{\mu\nu}I_4,\quad
g=\operatorname{diag}(1,-1,-1,-1)}.
\]

\Needspace{5\baselineskip}
\par\bigskip
\noindent\textbf{2.2.42}\quad Show that $\boldsymbol{\gamma}^{5}$, Eq.~(2.76),
anticommutes with all four $\boldsymbol{\gamma}^{\mu}$.
\solution
By definition, $\gamma^5=i\gamma^0\gamma^1\gamma^2\gamma^3$.
Use the tensor forms of 2.2.41 and the direct-product rule twice,
with $J^3=J^2J=-J$ and $\sigma_1\sigma_2\sigma_3=(i\sigma_3)\sigma_3=iI_2$
(Exercise 2.2.11):
\[
\gamma^1\gamma^2\gamma^3
=J^3\otimes(\sigma_1\sigma_2\sigma_3)
=(-J)\otimes(iI_2)=-iJ\otimes I_2.
\]
Hence, since $i(-i)=1$ and $\sigma_3J=\sigma_1$ (2.2.41),
\[
\gamma^5=i(\sigma_3\otimes I_2)(-iJ\otimes I_2)
=(\sigma_3J)\otimes I_2=\sigma_1\otimes I_2
=\begin{pmatrix}0&I_2\\I_2&0\end{pmatrix}.
\]
Now $\sigma_1\sigma_3=-\sigma_3\sigma_1$ (Exercise 2.2.11) and, with
$J=i\sigma_2$, $\sigma_1J=i\sigma_1\sigma_2=-i\sigma_2\sigma_1=-J\sigma_1$.
By the direct-product rule, therefore,
\[
\{\gamma^5,\gamma^0\}
=(\sigma_1\sigma_3+\sigma_3\sigma_1)\otimes I_2=0,
\]
\[
\{\gamma^5,\gamma^i\}
=(\sigma_1J+J\sigma_1)\otimes\sigma_i=0\quad(i=1,2,3).
\]
Thus $\gamma^5$ anticommutes with all four $\gamma^\mu$.

\Needspace{5\baselineskip}
\par\bigskip
\noindent\textbf{2.2.43}\quad In this problem, the summations are over
$\mu=0,1,2,3$. Define $g_{\mu\nu}=g^{\mu\nu}$ by the relations
\[
g_{00}=1;\quad g_{kk}=-1,\quad k=1,2,3;\quad g_{\mu\nu}=0,\quad\mu\neq\nu;
\]
and define $\gamma_{\mu}$ as
$\erratum{\sum_{\nu=0}^{3}g_{\nu\mu}\gamma^{\nu}}$. Using these
definitions, show that
\begin{enumerate}
\item[(a)] $\displaystyle\sum\gamma_{\mu}\gamma^{\alpha}\gamma^{\mu}
  =-2\gamma^{\alpha}$,
\item[(b)] $\displaystyle\sum\gamma_{\mu}\gamma^{\alpha}\gamma^{\beta}
  \gamma^{\mu}=4g^{\alpha\beta}\erratum{\mathbf1_4}$,
\item[(c)] $\displaystyle\sum\gamma_{\mu}\gamma^{\alpha}\gamma^{\beta}
  \gamma^{\nu}\gamma^{\mu}=-2\gamma^{\nu}\gamma^{\beta}\gamma^{\alpha}$.
\end{enumerate}
\solution
\erratum{The textbook prints the lowering sum without a summation
index and with $\gamma^\mu$ in the summand, although $\mu$ is the
free index. The intended definition is the contraction over the
dummy index $\nu$ shown above; for this diagonal metric it gives
$\gamma_\mu=g_{\mu\mu}\gamma^\mu$ with no sum on $\mu$. In (b), a
scalar denotes that scalar times $I_4$, made explicit in the
statement.}
The sums in (a)--(c) run over $\mu=0,1,2,3$.
Lowering an index gives $\gamma_0=\gamma^0$, $\gamma_i=-\gamma^i$.
The Clifford relation $\{\gamma^\mu,\gamma^\alpha\}=2g^{\mu\alpha}I_4$
of 2.2.41 implies
\[
\{\gamma_\mu,\gamma^\alpha\}=g_{\mu\mu}\{\gamma^\mu,\gamma^\alpha\}
=2g_{\mu\mu}g^{\mu\alpha}I_4=2\delta_\mu^\alpha I_4,\qquad
\sum_\mu\gamma_\mu\gamma^\mu=4I_4,
\]
because $g^{\mu\alpha}=0$ for $\alpha\ne\mu$ and
$g_{\mu\mu}g^{\mu\mu}=(\pm1)^2=1$.
For the second identity, the time term is $(\gamma^0)^2=I_4$ and each
spatial term is $-(\gamma^i)^2=I_4$.

\textbf{(a)} Move $\gamma_\mu$ past the first factor with
$\gamma_\mu\gamma^\alpha=2\delta_\mu^\alpha I_4-\gamma^\alpha\gamma_\mu$;
in the first term the Kronecker delta keeps only $\mu=\alpha$:
\[
\begin{aligned}
\sum_\mu\gamma_\mu\gamma^\alpha\gamma^\mu
&=\sum_\mu(2\delta_\mu^\alpha I_4-\gamma^\alpha\gamma_\mu)\gamma^\mu
=2\gamma^\alpha-\gamma^\alpha\sum_\mu\gamma_\mu\gamma^\mu\\
&=2\gamma^\alpha-\gamma^\alpha(4I_4)
=\boxed{-2\gamma^\alpha}.
\end{aligned}
\]
\textbf{(b)} Use the same step and then (a) with $\alpha$ replaced
by $\beta$:
\[
\begin{aligned}
\sum_\mu\gamma_\mu\gamma^\alpha\gamma^\beta\gamma^\mu
&=\sum_\mu(2\delta_\mu^\alpha I_4-\gamma^\alpha\gamma_\mu)
\gamma^\beta\gamma^\mu\\
&=2\gamma^\beta\gamma^\alpha
-\gamma^\alpha\sum_\mu\gamma_\mu\gamma^\beta\gamma^\mu\\
&=2\gamma^\beta\gamma^\alpha+2\gamma^\alpha\gamma^\beta
=2\{\gamma^\alpha,\gamma^\beta\}
=\boxed{4g^{\alpha\beta}I_4}.
\end{aligned}
\]
\textbf{(c)} Use the same step, then (b) with $(\alpha,\beta)$
replaced by $(\beta,\nu)$, and finally the Clifford relation in the
form $\gamma^\beta\gamma^\nu-2g^{\beta\nu}I_4=-\gamma^\nu\gamma^\beta$:
\[
\begin{aligned}
\sum_\mu\gamma_\mu\gamma^\alpha\gamma^\beta\gamma^\nu\gamma^\mu
&=2\gamma^\beta\gamma^\nu\gamma^\alpha
-\gamma^\alpha\sum_\mu\gamma_\mu\gamma^\beta\gamma^\nu\gamma^\mu\\
&=2\gamma^\beta\gamma^\nu\gamma^\alpha-4g^{\beta\nu}\gamma^\alpha\\
&=2(\gamma^\beta\gamma^\nu-2g^{\beta\nu}I_4)\gamma^\alpha
=\boxed{-2\gamma^\nu\gamma^\beta\gamma^\alpha}.
\end{aligned}
\]
These derivations include coincident free indices; no distinctness
assumption was made.

\Needspace{5\baselineskip}
\par\bigskip
\noindent\textbf{2.2.44}\quad If
$\mathsf{M}=\frac{1}{2}(1+\boldsymbol{\gamma}^{5})$, where
$\boldsymbol{\gamma}^{5}$ is given in Eq.~(2.76), show that
\[
\mathsf{M}^{2}=\mathsf{M}.
\]
\erratum{The textbook adds that this equation still holds when
$\boldsymbol{\gamma}^5$ is replaced by any other Dirac matrix listed in
Eq.~(2.76). That is true precisely for the replacements $\Gamma$ with
$\Gamma^2=\mathbf1_4$, not for every Dirac matrix (see the solution).}
\solution
Since $\gamma^5=\sigma_1\otimes I_2$ (Exercise 2.2.42),
$(\gamma^5)^2=\sigma_1^2\otimes I_2=I_4$. Expanding the square, in
which $I_4$ commutes with $\gamma^5$,
\[
M^2=\tfrac14(I_4+2\gamma^5+(\gamma^5)^2)
=\tfrac14(2I_4+2\gamma^5)=\tfrac12(I_4+\gamma^5)=M.
\]
For the note in the statement, let $\Gamma$ be any $4\times4$ matrix
and $M_\Gamma=(I_4+\Gamma)/2$. The same expansion gives
$M_\Gamma^2=\tfrac14(I_4+2\Gamma+\Gamma^2)$, so
\[
M_\Gamma^2-M_\Gamma
=\frac{I_4+2\Gamma+\Gamma^2}{4}-\frac{2I_4+2\Gamma}{4}
=\frac{\Gamma^2-I_4}{4},
\]
and $M_\Gamma^2=M_\Gamma$ exactly when $\Gamma^2=I_4$.
Each of the 16 Dirac matrices squares to $\pm I_4$:
\[
\begin{array}{c|c}
\Gamma&\Gamma^2\\\hline
I_4,\ \gamma^5,\ \gamma^0,\ \gamma^5\gamma^i,\
\sigma^{ij}\ (1\le i<j\le3)&+I_4\\
\gamma^i,\ \gamma^5\gamma^0,\ \sigma^{0i}\ (i=1,2,3)&-I_4
\end{array}
\]
where $\sigma^{\mu\nu}=i\gamma^\mu\gamma^\nu$ for $\mu<\nu$.
The entries follow from $(\gamma^0)^2=(\gamma^5)^2=I_4$,
$(\gamma^i)^2=-I_4$, and the anticommutation relations of
Exercises 2.2.41 and 2.2.42. For example, moving the second
$\gamma^5$ to the left past $\gamma^\mu$,
$(\gamma^5\gamma^\mu)^2=\gamma^5\gamma^\mu\gamma^5\gamma^\mu
=-(\gamma^5)^2(\gamma^\mu)^2=-(\gamma^\mu)^2$, and, for $\mu\ne\nu$,
$(i\gamma^\mu\gamma^\nu)^2=-\gamma^\mu\gamma^\nu\gamma^\mu\gamma^\nu
=\gamma^\mu\gamma^\mu\gamma^\nu\gamma^\nu=(\gamma^\mu)^2(\gamma^\nu)^2$.
Thus the replacement works for the matrices in the first row of the
table and fails for the seven in the second row, for which
$M_\Gamma^2-M_\Gamma=-I_4/2$.
Multiplying such a $\Gamma$ by $i$ makes its square $+I_4$ and
restores the projector construction.

\Needspace{5\baselineskip}
\par\bigskip
\noindent\textbf{2.2.45}\quad Prove that the 16 Dirac matrices form a linearly
independent set.
\solution
Let $\sigma_0=I_2$. All four $\sigma_a$ are Hermitian, so
$\sigma_a^\dagger\sigma_b=\sigma_a\sigma_b$, and the Pauli products of
Exercise 2.2.11 give
\[
\operatorname{tr}(\sigma_a^\dagger\sigma_b)=2\delta_{ab},
\qquad a,b=0,1,2,3.
\]
For $a=b$ the product is $\sigma_a^2=I_2$, of trace $2$; for distinct
indices the product is a traceless Pauli matrix or $\pm i$ times one.
Also, the diagonal entries of $X\otimes Y$ are the products
$X_{rr}Y_{ss}$, so
\[
\operatorname{tr}(X\otimes Y)=\sum_{r,s}X_{rr}Y_{ss}
=\operatorname{tr}X\,\operatorname{tr}Y.
\]
Therefore the 16 tensor products $T_{ab}=\sigma_a\otimes\sigma_b$ obey,
by $(X\otimes Y)^\dagger=X^\dagger\otimes Y^\dagger$ (Exercise 2.2.39)
and the direct-product rule,
\[
\operatorname{tr}(T_{ab}^\dagger T_{cd})
=\operatorname{tr}(\sigma_a^\dagger\sigma_c\otimes\sigma_b^\dagger\sigma_d)
=\operatorname{tr}(\sigma_a^\dagger\sigma_c)\,
\operatorname{tr}(\sigma_b^\dagger\sigma_d)
=4\delta_{ac}\delta_{bd}.
\]
The explicit correspondence with Eq.~(2.76) is
\[
\begin{array}{c|c}
\text{Dirac matrix}&\text{tensor form}\\\hline
I_4&\sigma_0\otimes\sigma_0\\
\gamma^5&\sigma_1\otimes\sigma_0\\
\gamma^0&\sigma_3\otimes\sigma_0\\
\gamma^5\gamma^0&-i\sigma_2\otimes\sigma_0\\
\gamma^i\ (i=1,2,3)&i\sigma_2\otimes\sigma_i\\
\gamma^5\gamma^i\ (i=1,2,3)&-\sigma_3\otimes\sigma_i\\
\sigma^{0i}\ (i=1,2,3)&i\sigma_1\otimes\sigma_i\\
\sigma^{12},\ \sigma^{13},\ \sigma^{23}
&\sigma_0\otimes\sigma_3,\ -\sigma_0\otimes\sigma_2,\
\sigma_0\otimes\sigma_1
\end{array}
\]
The first three rows are $I_4$ and the forms of $\gamma^5$ and
$\gamma^0$ found in Exercises 2.2.41 and 2.2.42, and the row of
$\gamma^i$ is $\gamma^i=J\otimes\sigma_i$ with $J=i\sigma_2$. The
others follow from the direct-product rule and the Pauli products of
Exercise 2.2.11:
\[
\begin{aligned}
\gamma^5\gamma^0&=\sigma_1\sigma_3\otimes\sigma_0
=-i\sigma_2\otimes\sigma_0,\qquad
\gamma^5\gamma^i=\sigma_1(i\sigma_2)\otimes\sigma_i
=-\sigma_3\otimes\sigma_i,\\
\sigma^{0i}&=i\gamma^0\gamma^i=i\,\sigma_3(i\sigma_2)\otimes\sigma_i
=-\sigma_3\sigma_2\otimes\sigma_i=i\sigma_1\otimes\sigma_i,\\
\sigma^{ij}&=i\gamma^i\gamma^j=i(i\sigma_2)^2\otimes\sigma_i\sigma_j
=-iI_2\otimes\sigma_i\sigma_j=\varepsilon_{ijk}I_2\otimes\sigma_k,
\end{aligned}
\]
where in the last line $i<j$, $k$ is the remaining index, and
$\sigma_i\sigma_j=i\varepsilon_{ijk}\sigma_k$.
The table uses each pair $(a,b)$ exactly once, so each Dirac matrix
is $\Gamma_r=\eta_rT_{a_rb_r}$ with a different pair $(a_r,b_r)$ and
a coefficient $\eta_r\in\{\pm1,\pm i\}$ of modulus $1$. Hence
\[
\operatorname{tr}(\Gamma_r^\dagger\Gamma_s)
=\eta_r^*\eta_s\operatorname{tr}(T_{a_rb_r}^\dagger T_{a_sb_s})
=4|\eta_r|^2\delta_{rs}=4\delta_{rs}.
\]
If $\sum_sc_s\Gamma_s=0$, multiply by $\Gamma_r^\dagger$ and
take the trace: $4c_r=0$ for every $r$. This proves linear
independence over $\mathbb C$, and hence over $\mathbb R$.
Other sign or phase conventions (for example $\gamma^0\gamma^5
=-\gamma^5\gamma^0$ in place of $\gamma^5\gamma^0$) only change some
$\eta_r$ and do not affect the argument.
The complex $4\times4$ matrix space has dimension $16$, so these
matrices also span it.

\Needspace{5\baselineskip}
\par\bigskip
\noindent\textbf{2.2.46}\quad If we assume that a given $4\times4$ matrix
$\mathsf{A}$ (with constant elements) can be written as a linear combination
of the 16 Dirac matrices (which we denote here as $\Gamma_{i}$)
\[
\mathsf{A}=\sum_{i=1}^{16}c_{i}\Gamma_{i},
\]
show that
\[
c_{i}\sim\mathrm{trace}(\mathsf{A}\Gamma_{i}).
\]
\solution
To avoid confusion with the imaginary unit, label the Dirac matrices
$\Gamma_r$, $r=1,\ldots,16$, and the coefficients $c_r$.
Write $\Gamma_r^2=s_rI_4$, where $s_r=\pm1$ as in the table of 2.2.44.
In 2.2.45 each $\Gamma_r=\eta_rT_r$, where $T_r$ is one of the
products $\sigma_a\otimes\sigma_b$ and $\eta_r\in\{\pm1,\pm i\}$.
Each $T_r$ is Hermitian, and $T_r^2=\sigma_a^2\otimes\sigma_b^2=I_4$;
hence $\Gamma_r^2=\eta_r^2I_4$, so $s_r=\eta_r^2$, and
\[
\Gamma_r^\dagger=\eta_r^*T_r=\frac{\eta_r^*}{\eta_r}\Gamma_r
=(\eta_r^*)^2\Gamma_r=s_r\Gamma_r,
\]
because $\eta_r^*/\eta_r=(\eta_r^*)^2/|\eta_r|^2$, $|\eta_r|=1$, and
$(\eta_r^*)^2=(\eta_r^2)^*=s_r$ is real. Since $s_r^2=1$, also
$\Gamma_r=s_r\Gamma_r^\dagger$, and the trace relation of 2.2.45 gives
\[
\operatorname{tr}(\Gamma_r\Gamma_j)
=s_r\operatorname{tr}(\Gamma_r^\dagger\Gamma_j)=4s_r\delta_{rj}.
\]
Multiplying $A=\sum_jc_j\Gamma_j$ by $\Gamma_r$ and taking the trace,
which is linear and satisfies
$\operatorname{tr}(\Gamma_j\Gamma_r)=\operatorname{tr}(\Gamma_r\Gamma_j)$,
\[
\operatorname{tr}(A\Gamma_r)=
\sum_jc_j\operatorname{tr}(\Gamma_j\Gamma_r)=4s_rc_r.
\]
Thus the full proportionality, with its normalization and sign, is
\[
\boxed{c_r=\frac1{4s_r}\operatorname{tr}(A\Gamma_r)
=\frac{s_r}{4}\operatorname{tr}(A\Gamma_r)
=\frac14\operatorname{tr}(\Gamma_r^\dagger A)}.
\]
The last form follows from $s_r\Gamma_r=\Gamma_r^\dagger$ and
$\operatorname{tr}(A\Gamma_r^\dagger)=\operatorname{tr}(\Gamma_r^\dagger A)$.
The sign $s_r$ matters for the basis matrices squaring to $-I_4$.

\Needspace{5\baselineskip}
\par\bigskip
\noindent\textbf{2.2.47}\quad The matrix
$\mathsf{C}=i\boldsymbol{\gamma}^{2}\boldsymbol{\gamma}^{0}$ is sometimes
called the charge conjugation matrix. Show that
$\mathsf{C}\boldsymbol{\gamma}^{\mu}\mathsf{C}^{-1}
=-(\boldsymbol{\gamma}^{\mu})^{T}$.
\solution
Use the specified Dirac representation
$\gamma^0=\sigma_3\otimes I_2$, $\gamma^i=i\sigma_2\otimes\sigma_i$
(Exercise 2.2.41). By the direct-product rule and
$\sigma_2\sigma_3=i\sigma_1$,
\[
C=i\gamma^2\gamma^0=i\,(i\sigma_2\sigma_3\otimes\sigma_2)
=-\sigma_2\sigma_3\otimes\sigma_2
=-i\sigma_1\otimes\sigma_2,
\]
so $C^2=(-i)^2\sigma_1^2\otimes\sigma_2^2=-I_4$ and $C^{-1}=-C$.
The transpose of a direct product is $(A\otimes B)^T=A^T\otimes B^T$,
since in the notation of 2.2.39
$((A\otimes B)^T)_{(k,a),(l,b)}=A_{lk}B_{ba}=(A^T)_{kl}(B^T)_{ab}$.
The Pauli transposes are
$\sigma_1^T=\sigma_1$, $\sigma_2^T=-\sigma_2$, $\sigma_3^T=\sigma_3$,
so $(\gamma^0)^T=\sigma_3\otimes I_2=\gamma^0$ and
$(\gamma^i)^T=(i\sigma_2)^T\otimes\sigma_i^T=-i\sigma_2\otimes\sigma_i^T$,
giving
\[
(\gamma^0)^T=\gamma^0,\quad(\gamma^2)^T=\gamma^2,\quad
(\gamma^1)^T=-\gamma^1,\quad(\gamma^3)^T=-\gamma^3.
\]
On the other side, the product $\gamma^2\gamma^0$ anticommutes with
$\gamma^0$ and $\gamma^2$: by $(\gamma^0)^2=I_4$, $(\gamma^2)^2=-I_4$,
and $\gamma^0\gamma^2=-\gamma^2\gamma^0$ (Exercise 2.2.41),
\[
\begin{aligned}
C\gamma^0&=i\gamma^2(\gamma^0)^2=i\gamma^2,&
\gamma^0C&=i\gamma^0\gamma^2\gamma^0=-i\gamma^2(\gamma^0)^2=-i\gamma^2,\\
C\gamma^2&=i\gamma^2\gamma^0\gamma^2=-i(\gamma^2)^2\gamma^0=i\gamma^0,&
\gamma^2C&=i(\gamma^2)^2\gamma^0=-i\gamma^0.
\end{aligned}
\]
It commutes with $\gamma^1,\gamma^3$, since moving either to the left
past the two factors $\gamma^0$ and $\gamma^2$ gives two minus signs:
$C\gamma^\mu=i\gamma^2\gamma^0\gamma^\mu=i\gamma^\mu\gamma^2\gamma^0
=\gamma^\mu C$ for $\mu=1,3$. Multiplying $C\gamma^\mu=\mp\gamma^\mu C$
on the right by $C^{-1}$ gives
\[
C\gamma^\mu C^{-1}=
\begin{cases}-\gamma^\mu,&\mu=0,2,\\+\gamma^\mu,&\mu=1,3.\end{cases}
\]
Comparison with the transposes proves
$\boxed{C\gamma^\mu C^{-1}=-(\gamma^\mu)^T}$ for every $\mu$.

\Needspace{5\baselineskip}
\par\bigskip
\noindent\textbf{2.2.48}\quad
\begin{enumerate}
\item[(a)] Show that, by substitution of the definitions of the
  $\boldsymbol{\gamma}^{\mu}$ matrices from Eqs.~(2.70) and (2.72), that the
  Dirac equation, Eq.~(2.73), takes the following form when written as
  $2\times2$ blocks (with $\psi_{L}$ and $\psi_{S}$ column vectors of
  dimension 2). Here $L$ and $S$ stand, respectively, for ``large'' and
  ``small'' because of their relative size in the nonrelativistic limit):
  \[
  \begin{pmatrix}
  mc^{2}-E&c(\sigma_{1}p_{1}+\sigma_{2}p_{2}+\sigma_{3}p_{3})\\
  \erratum{+}c(\sigma_{1}p_{1}+\sigma_{2}p_{2}+\sigma_{3}p_{3})&-mc^{2}-E
  \end{pmatrix}
  \begin{pmatrix}\psi_{L}\\\psi_{S}\end{pmatrix}=0.
  \]
\item[(b)] \erratum{Take $m>0$ and the positive-energy branch.}
  To reach the nonrelativistic limit, make the substitution
  $E=mc^{2}+\varepsilon$ and approximate $-2mc^{2}-\varepsilon$ by $-2mc^{2}$.
  Then write the matrix equation as two simultaneous two-component equations
  and show that they can be rearranged to yield
  \[
  \frac{1}{2m}\left(p_{1}^{2}+p_{2}^{2}+p_{3}^{2}\right)\psi_{L}
  =\varepsilon\psi_{L},
  \]
  which is just the Schr\"odinger equation for a free particle.
\item[(c)] Explain why is it reasonable to call $\psi_{L}$ and $\psi_{S}$
  ``large'' and ``small.''
\end{enumerate}
\solution
\erratum{The lower-left block in the printed question (textbook
p.~121) has an incorrect minus sign. Equation (2.73) requires the
positive sign shown in the corrected question above. The printed sign
would give negative free-particle kinetic energy upon elimination of
the small component: its second row would imply
$\psi_S\simeq-P\psi_L/(2mc)$, and its first would then give
$\varepsilon\psi_L=-\mathbf p^2\psi_L/(2m)$, where
$P=\boldsymbol{\sigma}\cdot\mathbf p$.}

\erratum{In the Instructor's Manual, the block Hamiltonian also omits
$c$ from its momentum terms. Equation (2.73) has $c\alpha_i p_i$;
all factors of $c$ are retained below.}

\textbf{(a)} Put $P=\boldsymbol{\sigma}\cdot\mathbf p
=\sigma_1p_1+\sigma_2p_2+\sigma_3p_3$. The definitions
in Eqs.~(2.70)--(2.73) give
\[
\gamma^0=\begin{pmatrix}I_2&0\\0&-I_2\end{pmatrix},\quad
\gamma^i=\begin{pmatrix}0&\sigma_i\\-\sigma_i&0\end{pmatrix},\quad
\alpha_i=\gamma^0\gamma^i=
\begin{pmatrix}0&\sigma_i\\\sigma_i&0\end{pmatrix};
\]
in the block product $\gamma^0\gamma^i$ the upper-right block is
$I_2\sigma_i=\sigma_i$, the lower-left block is
$(-I_2)(-\sigma_i)=\sigma_i$, and the diagonal blocks vanish. Hence
\[
mc^2\gamma^0+c\sum_i\alpha_ip_i-EI_4
=\begin{pmatrix}(mc^2-E)I_2&cP\\cP&(-mc^2-E)I_2\end{pmatrix},
\]
and with $\Psi=(\psi_L,\psi_S)^T$ the Dirac equation
$(mc^2\gamma^0+c\sum_i\alpha_ip_i)\Psi=E\Psi$ is
\[
\boxed{\begin{pmatrix}
(mc^2-E)I_2&cP\\cP&(-mc^2-E)I_2
\end{pmatrix}\begin{pmatrix}\psi_L\\\psi_S\end{pmatrix}=0}.
\]
Both off-diagonal blocks are the same.

\textbf{(b)} On the positive-energy branch set $E=mc^2+\varepsilon$
with $|\varepsilon|\ll mc^2$. Then $mc^2-E=-\varepsilon$ and
$-mc^2-E=-(2mc^2+\varepsilon)$, so the two block rows of (a) are the
two-component equations
\[
-\varepsilon\psi_L+cP\psi_S=0,\qquad
cP\psi_L-(2mc^2+\varepsilon)\psi_S=0.
\]
The second gives $\psi_S=cP\psi_L/(2mc^2+\varepsilon)$.
Substitution into the first yields, since $P^2=\mathbf p^2I_2$ with
$\mathbf p^2=p_1^2+p_2^2+p_3^2$ (Exercise 2.2.40),
\[
\varepsilon\psi_L=\frac{c^2P^2}{2mc^2+\varepsilon}\psi_L
=\frac{c^2\mathbf p^2}{2mc^2+\varepsilon}\psi_L.
\]
Approximating $-2mc^2-\varepsilon$ by $-2mc^2$, as the problem
directs, replaces $(2mc^2+\varepsilon)^{-1}$ by $(2mc^2)^{-1}$; since
$c^2/(2mc^2)=1/(2m)$, this gives
\[
\boxed{\frac{\mathbf p^2}{2m}\psi_L=\varepsilon\psi_L}.
\]
With $p_i=-i\hbar\partial_i$, this is
$-\hbar^2\nabla^2\psi_L/(2m)=\varepsilon\psi_L$.
As an exact check before approximation, take a plane wave of
momentum $\mathbf p$ (so that $\mathbf p$ is a numerical vector)
with $\psi_L\ne0$. Multiplying by $2mc^2+\varepsilon$ gives
$\varepsilon(2mc^2+\varepsilon)=c^2\mathbf p^2$; since
$\varepsilon(2mc^2+\varepsilon)=(E-mc^2)(E+mc^2)=E^2-m^2c^4$, this is
$E^2=m^2c^4+c^2\mathbf p^2$.

\textbf{(c)} For such a state $2mc^2+\varepsilon=E+mc^2>0$, and the
exact relation of (b) is $\psi_S=cP\psi_L/(E+mc^2)$. For real
$\mathbf p$, $P$ is Hermitian, so $P^\dagger P=P^2=\mathbf p^2I_2$ and
\[
\|\psi_S\|^2=\frac{c^2\,\psi_L^\dagger P^\dagger P\psi_L}{(E+mc^2)^2}
=\frac{c^2\mathbf p^2}{(E+mc^2)^2}\,\|\psi_L\|^2.
\]
Hence, with $E\simeq mc^2$ in the nonrelativistic regime,
\[
\frac{\|\psi_S\|}{\|\psi_L\|}
=\frac{c|\mathbf p|}{E+mc^2}
\simeq\frac{|\mathbf p|}{2mc}\ll1
\quad\text{when }|\mathbf p|\ll mc.
\]
This explains the names on the positive-energy, nonrelativistic
branch. At rest the small component vanishes; the ratio need not be
small in the relativistic regime, and it tends to $1$ as
$|\mathbf p|\to\infty$, where $E\simeq c|\mathbf p|$.

\Needspace{5\baselineskip}
\par\bigskip
\noindent\textbf{2.2.49}\quad Show that it is consistent with the requirements
that they must satisfy to take the Dirac gamma matrices to be (in $2\times2$
block form)
\[
\boldsymbol{\gamma}^{0}=\begin{pmatrix}0&\mathbf{1}_{2}\\
\mathbf{1}_{2}&0\end{pmatrix},\quad
\boldsymbol{\gamma}^{i}=\begin{pmatrix}0&\sigma_{i}\\
-\sigma_{i}&0\end{pmatrix},\ (i=1,2,3).
\]
This choice for the gamma matrices is called the Weyl representation.
\solution
Write the proposed Weyl matrices as
\[
\gamma^0=\sigma_1\otimes I_2,\qquad
\gamma^i=J\otimes\sigma_i,\quad
J=i\sigma_2=\begin{pmatrix}0&1\\-1&0\end{pmatrix};
\]
the $\gamma^i$ are the same as in the Dirac representation of 2.2.41.
Then $(\gamma^0)^2=\sigma_1^2\otimes I_2=I_4$ and
$(\gamma^i)^2=J^2\otimes\sigma_i^2=-I_4$.
By the direct-product rule $\gamma^0\gamma^i=\sigma_1J\otimes\sigma_i$
and $\gamma^i\gamma^0=J\sigma_1\otimes\sigma_i$, where
\[
\sigma_1J=\begin{pmatrix}0&1\\1&0\end{pmatrix}
\begin{pmatrix}0&1\\-1&0\end{pmatrix}
=\begin{pmatrix}-1&0\\0&1\end{pmatrix}=-\sigma_3,\qquad
J\sigma_1=\begin{pmatrix}1&0\\0&-1\end{pmatrix}=\sigma_3.
\]
Hence
\[
\{\gamma^0,\gamma^i\}=(\sigma_1J+J\sigma_1)\otimes\sigma_i=0.
\]
For spatial indices, as in 2.2.41,
$\{\gamma^i,\gamma^j\}=J^2\otimes\{\sigma_i,\sigma_j\}
=-2\delta_{ij}I_4$. All the required relations (2.74)--(2.75) hold.
Also $(\gamma^0)^\dagger=\sigma_1^\dagger\otimes I_2=\gamma^0$ and,
since $J^\dagger=J^T=-J$ and $\sigma_i^\dagger=\sigma_i$,
$(\gamma^i)^\dagger=J^\dagger\otimes\sigma_i^\dagger=-\gamma^i$. With
the anticommutation just proved,
\[
\alpha_i^\dagger=(\gamma^0\gamma^i)^\dagger
=(\gamma^i)^\dagger(\gamma^0)^\dagger
=-\gamma^i\gamma^0=\gamma^0\gamma^i=\alpha_i.
\]
The free Dirac Hamiltonian $mc^2\gamma^0+c\sum_i\alpha_ip_i$
therefore has the required Hermiticity as well.

\Needspace{5\baselineskip}
\par\bigskip
\noindent\textbf{2.2.50}\quad Show that the Dirac equation separates into
independent $2\times2$ blocks in the Weyl representation (see Exercise 2.2.49)
in the limit that the mass $m$ approaches zero. This observation is important
in the ultra relativistic regime where the rest mass is inconsequential, or
for particles of negligible mass (e.g., neutrinos).
\solution
\erratum{The Instructor's Manual's momentum blocks again lack their
factors of $c$. With the units of Eq.~(2.73), the decoupled equations
contain $\pm c\boldsymbol{\sigma}\cdot\mathbf p$.}
In the Weyl representation, block multiplication gives
\[
\alpha_i=\gamma^0\gamma^i=
\begin{pmatrix}0&I_2\\I_2&0\end{pmatrix}
\begin{pmatrix}0&\sigma_i\\-\sigma_i&0\end{pmatrix}
=\begin{pmatrix}-\sigma_i&0\\0&\sigma_i\end{pmatrix},
\qquad
mc^2\gamma^0=\begin{pmatrix}0&mc^2I_2\\mc^2I_2&0\end{pmatrix},
\]
so $c\sum_i\alpha_ip_i$ is block diagonal with diagonal blocks
$-c\boldsymbol{\sigma}\cdot\mathbf p$ and
$c\boldsymbol{\sigma}\cdot\mathbf p$.
With $\Psi=(\psi_1,\psi_2)^T$, Eq.~(2.73),
$(mc^2\gamma^0+c\sum_i\alpha_ip_i)\Psi=E\Psi$, is
\[
\begin{pmatrix}
-c\boldsymbol{\sigma}\cdot\mathbf p&mc^2I_2\\
mc^2I_2&c\boldsymbol{\sigma}\cdot\mathbf p
\end{pmatrix}
\begin{pmatrix}\psi_1\\\psi_2\end{pmatrix}
=E\begin{pmatrix}\psi_1\\\psi_2\end{pmatrix}.
\]
Equivalently,
\[
-c\boldsymbol{\sigma}\cdot\mathbf p\,\psi_1+mc^2\psi_2=E\psi_1,\qquad
mc^2\psi_1+c\boldsymbol{\sigma}\cdot\mathbf p\,\psi_2=E\psi_2.
\]
At $m=0$ the coupling terms disappear:
\[
\boxed{-c\boldsymbol{\sigma}\cdot\mathbf p\,\psi_1=E\psi_1,\qquad
c\boldsymbol{\sigma}\cdot\mathbf p\,\psi_2=E\psi_2}.
\]
Each two-component equation can be solved independently. Squaring
either gives the massless relation: applying
$-c\boldsymbol{\sigma}\cdot\mathbf p$ to both sides of the first and
using $(\boldsymbol{\sigma}\cdot\mathbf p)^2=\mathbf p^2I_2$
(Exercise 2.2.40),
\[
c^2\mathbf p^2\psi_1=E(-c\boldsymbol{\sigma}\cdot\mathbf p\,\psi_1)
=E^2\psi_1,
\]
and likewise $c^2\mathbf p^2\psi_2=E^2\psi_2$. For a plane wave with a
nonzero component this is $E^2=c^2\mathbf p^2$.

\Needspace{5\baselineskip}
\par\bigskip
\noindent\textbf{2.2.51}\quad
\begin{enumerate}
\item[(a)] Given $\mathbf{r}'=\mathsf{U}\mathbf{r}$, with $\mathsf{U}$ a
  unitary matrix and $\mathbf{r}$ a (column) vector with complex elements,
  show that the magnitude of $\mathbf{r}$ is invariant under this operation.
\item[(b)] The \erratum{square} matrix $\mathsf{U}$ transforms any column vector $\mathbf{r}$
  with complex elements into $\mathbf{r}'$, leaving the magnitude invariant:
  $\mathbf{r}^{\dagger}\mathbf{r}=\mathbf{r}'^{\dagger}\mathbf{r}'$. Show that
  $\mathsf{U}$ is unitary.
\end{enumerate}
\solution
\textbf{(a)} For a complex vector,
$\|\mathbf r\|^2=\mathbf r^\dagger\mathbf r$.
If $\mathbf r'=U\mathbf r$ with $U^\dagger U=I$, then, since
$(U\mathbf r)^\dagger=\mathbf r^\dagger U^\dagger$ (Exercise 2.2.32),
\[
\|\mathbf r'\|^2=(U\mathbf r)^\dagger U\mathbf r
=\mathbf r^\dagger U^\dagger U\mathbf r
=\mathbf r^\dagger\mathbf r.
\]
Taking nonnegative square roots proves invariance of the magnitude.

\textbf{(b)} \erratum{In the Instructor's Manual,
$\mathbf r'^{\dagger}\mathbf r$ must be $\mathbf r'^{\dagger}\mathbf r'$.
Here $U$ is square and acts on vectors of the
same dimension. For rectangular $U$, norm preservation implies only
$U^\dagger U=I$, not two-sided unitarity.}
Set $H=U^\dagger U-I$. Since
$\mathbf r'^\dagger\mathbf r'=\mathbf r^\dagger U^\dagger U\mathbf r$
as in (a), the hypothesis is
$\mathbf r^\dagger H\mathbf r=0$ for every complex $\mathbf r$;
also $H^\dagger=(U^\dagger U)^\dagger-I=H$, so $H_{ji}=H_{ij}^*$.
Let $e_i$ be the $i$th unit column vector, so that
$e_i^\dagger He_j=H_{ij}$.
Choose $\mathbf r=e_i$ to get $H_{ii}=0$.
For $i\ne j$, choose $\mathbf r=e_i+e_j$:
\[
0=(e_i+e_j)^\dagger H(e_i+e_j)=H_{ii}+H_{ij}+H_{ji}+H_{jj}
=H_{ij}+H_{ij}^*=2\operatorname{Re}H_{ij}.
\]
Next choose $\mathbf r=e_i+ie_j$, for which
$\mathbf r^\dagger=e_i^\dagger-ie_j^\dagger$ and $(-i)(i)=1$:
\[
0=H_{ii}+iH_{ij}-iH_{ji}+H_{jj}
=i(H_{ij}-H_{ij}^*)=i\cdot2i\operatorname{Im}H_{ij}
=-2\operatorname{Im}H_{ij}.
\]
Every entry vanishes, so $H=0$ and $U^\dagger U=I$.
Since $U$ is square, taking determinants yields
$\det U^\dagger\det U=(\det U)^*\det U=|\det U|^2=1$.
Hence $\det U\ne0$ and $U$ is invertible; multiplying $U^\dagger U=I$
on the right by $U^{-1}$ gives $U^\dagger=U^{-1}$, and then
$UU^\dagger=UU^{-1}=I$.
Thus $U$ is unitary.
